%%%%%%%%%%%%%%%%%%%%%%%%%%%%%%%%%%%%%%%%%%%%%%%%%%%%%%%%%%%%%%%%%%%%%%
% Lớp 11 | Hàm số mũ và hàm số lôgarít | Phép tính luỹ thừa
% Dạng 1D6?1-1: Tính giá trị của biểu thức chứa lũy thừa
% Số câu: 126  (sinh tự động bởi tach_cau.py; từ nay thầy sửa trực tiếp file này)
%%%%%%%%%%%%%%%%%%%%%%%%%%%%%%%%%%%%%%%%%%%%%%%%%%%%%%%%%%%%%%%%%%%%%%

%%%%%%%%%%%%%%%% TRẮC NGHIỆM 4 PHƯƠNG ÁN %%%%%%%%%%%%%%%%

% nguồn: Ok_Toan11_TracNghiem_Lan2
\begin{ex}%[1D6N1-1]%[ma:5d2c88c3]
	 Biểu thức $P=\sqrt[5]{-4}\cdot\sqrt[5]{8}$ có giá trị bằng
	\choice 
		{$4\sqrt{2}$} 
		 {\True $-2$}
		 {$2$}
		 {$-4\sqrt{2}$}
	\loigiai{ 
		Ta có $P=\sqrt[5]{-4}\cdot\sqrt[5]{8}=\sqrt[5]{(-32)}=-2$. 
	} 
\end{ex}

% nguồn: Ok_Toan11_TracNghiem_Lan2
\begin{ex}%[1D6N1-1]%[ma:64a00070]
	Cho $a=\dfrac{1}{256}$ và $b=\dfrac{1}{27}$.  Tính $A=a^{-\tfrac{3}{4}}+b^{-\tfrac{4}{3}}$.
	\choice 
		{$23$} 
		 {$89$}
		{\True  $145$}
		 {$26$}
	\loigiai{ 
		Thay $a=\dfrac{1}{256}$, $b=\dfrac{1}{27}$ vào $A=a^{-\tfrac{3}{4}}+b^{-\tfrac{4}{3}}$  ta được
		$$A=\left(\dfrac{1}{256}\right)^{-\tfrac{3}{4}}+\left(\dfrac{1}{27}\right)^{-\tfrac{4}{3}} =\left(4^{-4}\right)^{-\tfrac{3}{4}}+\left(3^{-3}\right)^{-\tfrac{4}{3}} =4^3+3^4=145.$$ 
	} 
\end{ex}

% nguồn: Ok_Toan11_TracNghiem_Lan2
\begin{ex}%[1D6N1-1]%[ma:92d77bc2]
	Cho $a$ là số thực dương. Biểu thức $a^3\cdot \sqrt[3]{a^2}$ được viết dưới dạng lũy thừa với số mũ  hữu tỉ là
	\choice 
		{\True  $a^{\tfrac{11}{3}}$}
		 {$a^2$}
		{$a^{\tfrac{5}{3}}$} 
		{$a^{\tfrac{8}{3}}$} 
	 \loigiai{
	 Ta có $a^3\cdot \sqrt[3]{a^2}=a^3\cdot a^{\tfrac{2}{3}}=a^{3+\tfrac{2}{3}}=a^{\tfrac{11}{3}}$.}
\end{ex}

% nguồn: Ok_Toan11_TracNghiem_Lan2
\begin{ex}%[1D6N1-1]%[ma:d86174f3]
	Giá trị $\sqrt[3]{2\,021}\cdot\sqrt[5]{2\,021}$ viết dưới dạng lũy  thữa với số mũ hữu tỷ là
	\choice 
		 {$2\,021^{\tfrac{2}{5}}$}
		 {$2\,021^{\tfrac{1}{15}}$}
		{\True  $2\,021^{\tfrac{8}{15}}$}
		{$2\,021^{\tfrac{1}{10}}$} 
	\loigiai{ 
		Ta có  $\sqrt[3]{2\,021}\cdot\sqrt[5]{2\,021}=2\,021^{\tfrac{1}{3}}\cdot 2\,021^{\tfrac{1}{5}}=2\,021^{\tfrac{8}{15}}$.
	 }
\end{ex}

% nguồn: Ok_Toan11_TracNghiem_Lan2
\begin{ex}%[1D6N1-1]%[ma:3541ecb6]
	Cho $a$  là số thực dương. Khẳng định nào sau đây đúng?
	 \choice
		{$a^m+a^n=a^{m+n}$} 
		 {$a^m\cdot a^n=a^{m-n}$}
		{\True  $(a^m)^n=(a^n)^m$}
		{$\dfrac{a^m}{a^n}=a^{n-m}$} 
	\loigiai{ 
	Tính chất  lũy thừa.}
\end{ex}

% nguồn: Ok_Toan11_TracNghiem_Lan2
\begin{ex}%[1D6N1-1]%[ma:78bef468]
	Tính giá trị của biểu thức $P=\left(7+4\sqrt{3}\right)^{2\, 017}\cdot \left(4\sqrt{3}-7\right)^{2\,  016}$.
	 \choice
		 {$P=\left(7+4\sqrt{3}\right)^{2\, 016}$}
		 {$P=1$}
		{$P=7-4\sqrt{3}$} 
		{\True  $P=7+4\sqrt{3}$}
	\loigiai{ 
		 Ta có
		\allowdisplaybreaks 
		\begin{eqnarray*}
			P  &= & \left(7+4\sqrt{3}\right)\cdot \left[\left(7+4\sqrt{3}\right)\cdot \left(4\sqrt{3}-7\right)\right]^{2\, 016}\\
			&= &  \left(7+4\sqrt{3}\right)\cdot \left(-1\right)^{2\, 016} =7+4\sqrt{3}.
		\end{eqnarray*}
	} 
\end{ex}

% nguồn: Ok_Toan11_TracNghiem_Lan2
\begin{ex}%[1D6N1-1]%[ma:3aafbef6]
	Viết biểu thức $P=\sqrt[3]{x}\cdot \sqrt[4]{x}$, $\left(x>0\right)$ dưới dạng lũy thừa với số mũ hữu  tỷ
	\choice 
		 {$P=x^{\tfrac{5}{4}}$}
		{$P=x^{\tfrac{1}{12}}$} 
		{$P=x^{\tfrac{1}{7}}$} 
		 {\True $P=x^{\tfrac{7}{12}}$}
	\loigiai{ 
	Ta có $P=\sqrt[3]{x}\cdot \sqrt[4]{x}=x^{\tfrac{1}{3}}\cdot  x^{\tfrac{1}{4}}=x^{\tfrac{1}{3}+\tfrac{1}{4}}=x^{\tfrac{7}{12}}$.}
\end{ex}

% nguồn: Ok_Toan11_TracNghiem_Lan2
\begin{ex}%[1D6N1-1]%[ma:185a85b8]
	Với $a$ là số thực  dương tùy ý, $a^4\cdot a^{\tfrac{1}{2}}$ bằng
	 \choice
		 {$a^8$}
		{$a^2$} 
		{$a^{\tfrac{7}{2}}$} 
		 {\True $a^{\tfrac{9}{2}}$}
	 \loigiai{
	Ta có $a^4\cdot  a^{\tfrac{1}{2}}=a^{4+\tfrac{1}{2}}=a^{\tfrac{9}{2}}$.}
\end{ex}

% nguồn: Ok_Toan11_TracNghiem_Lan2
\begin{ex}%[1D6N1-1]%[ma:5e038ecd]
	Cho biểu thức $P=\sqrt[3]{\dfrac{2}{3}\cdot\sqrt[3]{\dfrac{2}{3}\cdot\sqrt{\dfrac{2}{3}}}}$. Mệnh đề nào  trong các mệnh đề sau là đúng?
	\choice 
		 {$P=\left(\dfrac{2}{3}\right)^{\tfrac{1}{8}}$}
		{$P=\left(\dfrac{2}{3}\right)^{18}$} 
		 {$P=\left(\dfrac{2}{3}\right)^{\tfrac{1}{18}}$}
		{\True  $P=\left(\dfrac{2}{3}\right)^{\tfrac{1}{2}}$}
	 \loigiai{
		Ta có $P=\sqrt[3]{\dfrac{2}{3}\cdot\sqrt[3]{\left(\dfrac{2}{3}\right)^{\tfrac{3}{2}}}}  =\sqrt[3]{\left(\dfrac{2}{3}\right)^{\tfrac{3}{2}\cdot\tfrac{1}{3}+1}} =\sqrt[3]{\left(\dfrac{2}{3}\right)^{\tfrac{3}{2}}} =\left(\dfrac{2}{3}\right)^{\tfrac{1}{2}}$.
	} 
\end{ex}

% nguồn: Ok_Toan11_TracNghiem_Lan2
\begin{ex}%[1D6N1-1]%[ma:e4fcbbc6]
	Cho số thực dương $a$ và số nguyên dương $n$ tùy ý. Mệnh đề nào dưới đây  đúng?
	\choice 
		 {$\sqrt{a^n}=a^{2+n}$}
		{$\sqrt{a^n}=a^{2n}$} 
		 {$\sqrt{a^n}=a^{\tfrac{2}{n}}$}
		{\True $\sqrt{a^n}=a^{\tfrac{n}{2}}$} 
	\loigiai{ 
	Ta  có $\sqrt{a^n}=a^{\tfrac{n}{2}}$.}
\end{ex}

% nguồn: Ok_Toan11_TracNghiem_Lan2
\begin{ex}%[1D6N1-1]%[ma:d375f46e]
	Giá trị của  $27^{\tfrac{1}{3}}$ bằng
	\choice 
		{$6$} 
		 {$81$}
		 {$9$}
		 {\True $3$}
	\loigiai{ 
		Ta có  $27^{\tfrac{1}{3}}=\sqrt[3]{27}=3$.
	} 
\end{ex}

% nguồn: Ok_TracNghiemLop11_Lan1
\begin{ex}%[1D6N1-1]%[ma:67c340dd]
	Cho $a$  và $b$ là hai số thực dương thỏa mãn $a^2b^3=16$. Giá trị của $2\log_2a+3\log_2b$ bằng
	 \choice
		{$2$} 
		 {$8$}
		{$16$} 
		 {\True $4$}
	 \loigiai{
	Ta có $2\log_2a+3\log_2b=\log_2\left(a^2b^3\right)=\log_216=4$.} 
\end{ex}

% nguồn: Ok_TracNghiemLop11_Lan1
\begin{ex}%[1D6N1-1]%[ma:e489bb62]
	Cho $x,y>0$ và $\alpha, \beta \in \mathbb{R}$. Tìm đẳng thức \textbf{sai} dưới đây. 
	 \choice
		 {$(xy)^{\alpha}=x^{\alpha} \cdot y^{\alpha}$}
		{\True  $x^{\alpha}+y^{\alpha}=(x+y)^{\alpha}$}
		{$(x^{\alpha})^{\beta}=x^{\alpha \beta}$} 
		 {$x^{\alpha} \cdot x^{\beta}=x^{\alpha + \beta}$}
	 \loigiai{
		Theo tính chất của lũy thừa thì đẳng thức $x^{\alpha}+y^{\alpha}=(x+y)^{\alpha}$  \textbf{sai}.
	 }
\end{ex}

% nguồn: Ok_TracNghiemLop11_Lan1
\begin{ex}%[1D6N1-1]%[ma:2bb8f437]
	Thu gọn  biểu thức $A=\sqrt[4]{a^3} \cdot a$ với $a$ là số thực dương ta được
	\choice 
		 {$A=a^{\tfrac{5}{2}}$}
		{\True  $A=a^{\tfrac{7}{4}}$}
		 {$A=a^{\tfrac{3}{4}}$}
		 {$A=a^{\tfrac{1}{4}}$}
	\loigiai{ 
		Ta có $A=\sqrt[4]{a^3} \cdot  a=a^{\tfrac{3}{4}}\cdot a=a^{\tfrac{7}{4}}$.
	 }
\end{ex}

% nguồn: Ok_TracNghiemLop11_Lan1
\begin{ex}%[1D6N1-1]%[ma:d511cb4e]
	Với $a>0$, $b> 0$, $\alpha$, $\beta$ là các số  thực bất kì, đẳng thức nào sau đây \textbf{sai}?
	 \choice
		 {$a^{\alpha}\cdot a^{\beta}=a^{\alpha+\beta}$}
		{$a^{\alpha}\cdot b^{\alpha}=\left(ab\right)^{\alpha}$} 
		{$\dfrac{a^{\alpha}}{a^{\beta}}=a^{\alpha-\beta}$} 
		{\True $\dfrac{a^{\alpha}}{b^{\beta}}=\left(\dfrac{a}{b}\right)^{\alpha-\beta}$} 
	 \loigiai{
	Công thức $\dfrac{a^{\alpha}}{b^{\beta}}=\left(\dfrac{a}{b}\right)^{\alpha-\beta}$ là công thức sai.} 
\end{ex}

% nguồn: Ok_TracNghiemLop11_Lan1
\begin{ex}%[1D6N1-1]%[ma:4d45238a]
	 Cho các số thực $a,b,m,n\left(a,b>0\right)$. Khẳng định nào sau đây là \textbf{đúng}?
	 \choice
		 {$\dfrac{a^m}{a^n}=\sqrt{a^m}$}
		{$\left(a^m\right)^n=a^{m+n}$} 
		 {$\left(a+b\right)^m=a^m+b^m$}
		{\True $a^m\cdot a^n=a^{m+n}$} 
	\loigiai{ 
		Ta  có
		\begin{itemize}
			\item  $\dfrac{a^m}{a^n}=a^{m-n}\Rightarrow $ loại phương án $\dfrac{a^m}{a^n}=\sqrt{a^m}$.
			\item $\left(a^m\right)^n=a^{m\cdot  n}\Rightarrow $ loại phương án $\left(a^m\right)^n=a^{m+n}$.
			 \item $\left(1+1\right)^2\ne{1^2}+1^2\Rightarrow $ loại phương án $\left(a+b\right)^m=a^m+b^m$.
			\item  $a^m\cdot a^n=a^{m+n}\Rightarrow $ chọn phương án $a^m\cdot a^n=a^{m+n}$.
		\end{itemize}
	 }
\end{ex}

% nguồn: Ok_TracNghiemLop11_Lan1
\begin{ex}%[1D6N1-1]%[ma:fbccf455]
	Cho các số dương $a\ne 1$ và các số thực $\alpha $, $\beta $.  Đẳng thức nào sau đây là sai?
	\choice 
		 {$a^\alpha\cdot a^\beta=a^{\alpha+\beta}$}
		{\True $a^\alpha\cdot  a^\beta=a^{\alpha\beta}$}
		{$\dfrac{a^\alpha}{a^\beta}=a^{\alpha-\beta}$} 
		{$\left(a^\alpha\right)^\beta=a^{\alpha\beta}$} 
	 \loigiai{
	 }
\end{ex}

% nguồn: Ok_TracNghiemLop11_Lan1
\begin{ex}%[1D6N1-1]%[ma:35016822]
	Với $m$ là số  thực dương tùy ý, biểu thức $m^{\tfrac{10}{3}}:m^{\tfrac{1}{3}}$ bằng
	\choice 
		{\True  $m^3$}{$m^{10}$}{$m^{\tfrac{11}{3}}$}{$m^{\tfrac{10}{3}}$}
	 \loigiai{
		Ta có $m^{\tfrac{10}{3}}:  m^{\tfrac{1}{3}}=m^{\tfrac{10}{3}-\tfrac{1}{3}}=m^3$.
	 }
\end{ex}

% nguồn: Ok_TracNghiemLop11_Lan1
\begin{ex}%[1D6N1-1]%[ma:42d58f43]
	Cho số thực $a$ dương và số  hữu tì $r=\dfrac{m}{n}$, trong đó $m$ là một số nguyên và $n$ là số nguyên dương. Lũy thừa của $a$ với số mũ $r$, ký hiệu $a^r$, được xác định bời
	 \choice
		 {$a^r=a^{\frac{m}{n}}=\sqrt[m]{a^n}$}{$a^r=a^{\frac{m}{n}}=\sqrt[m n]{a}$}{$a^r=a^{\frac{m}{n}}=\sqrt[m+n]{a}$}{\True $r=a^{\frac{m}{n}}=\sqrt{a^m}$}
	 \loigiai{
		Ta có $r=a^{\frac{m}{n}}=\sqrt{a^m}$. 
	} 
\end{ex}

% nguồn: Ok_TracNghiemLop11_Lan1
\begin{ex}%[1D6N1-1]%[ma:7ec60664]
	Cho $a>0; \,a  \neq 1;\, \alpha, \beta \in \mathbb{R}$. Khẳng định nào sau đây \textbf{sai}?
	 \choice
		{$\dfrac{a^\alpha}{a^\beta}=a^{\alpha-\beta}$} 
		{$a^\alpha \cdot a^\beta=a^{\alpha+\beta}$} 
		{\True  $a^\alpha+a^\beta=a^{\alpha+\beta}$}
		 {$\left(a^\alpha\right)^\beta=a^{\alpha\cdot \beta}$}
	 \loigiai{
		Khẳng định  sai: $a^\alpha+a^\beta=a^{\alpha+\beta}$.
	} 
\end{ex}

% nguồn: Ok_TracNghiemLop11_Lan1
\begin{ex}%[1D6N1-1]%[ma:e4cdad1d]
	 Với $a$ là số thực dương tùy ý, biểu thức $a^{\frac{5}{3}}\cdot a^{\frac{1}{3}}$ là
	 \choice
		{$a^5$} 
		{$a^{\frac{^5}{9}}$} 
		 {$a^{\frac{4}{3}}$}
		{\True  $a^2$}
	\loigiai{ 
		 Ta có $a^{\frac{5}{3}}\cdot a^{\frac{1}{3}}=a^2$.
	} 
\end{ex}

% nguồn: Ok_TracNghiemLop11_Lan1
\begin{ex}%[1D6N1-1]%[ma:48dd26c8]
	Rút gọn biểu  thức $Q=b^{\frac{5}{3}}:\sqrt[3]b$ với $b>0$.
	 \choice
		 {$Q=b^{-\frac{4}{3}}$}
		 {\True $Q=b^{\frac{4}{3}}$}
		{$Q=b^{\frac{5}{9}}$} 
		{$Q=b^2$} 
	 \loigiai{
	$Q=b^{\frac{5}{3}}:\sqrt[3]b=b^{\frac{5}{3}}:b^{\frac{1}{3}}=b^{\frac{4}{3}}$.} 
\end{ex}

% nguồn: Ok_TracNghiemLop11_Lan1
\begin{ex}%[1D6N1-1]%[ma:9be06c01]
	Với $a$ là số thực dương  tùy ý, $\sqrt[3]{a^2}$ bằng
	\choice 
		{$a^{\frac{1}{6}}$} 
		 {$a^6$}
		 {\True $a^{\frac{2}{3}}$}
		{$a^{\frac{3}{2}}$} 
	\loigiai{ 
	Với mọi số thực dương $a$ ta có $\sqrt[3]{a^2}=a^{\frac{2}{3}}$.} 
\end{ex}

% nguồn: Ok_TracNghiemLop11_Lan1
\begin{ex}%[1D6N1-1]%[ma:29fed22a]
	Chọn khẳng định \textbf{sai}. 
	\choice 
		{Nếu $0<a<1$ thì $\log_a u>\log_a v \Leftrightarrow  0<u<v$}
		{Nếu $a>1$  thì $a^u >a^v \Leftrightarrow u>v$}
		{\True Nếu $a>1$ thì $\log_a u>\log_a v \Leftrightarrow u>v$} 
		{Nếu  $0<a<1$ thì $a^u >a^v \Leftrightarrow u<v$}
	\loigiai{Theo tính chất của  các hàm số hàm số logarit, nếu $a>1$ thì $\log_a u>\log_a v \Leftrightarrow u>v>0$.}
\end{ex}

% nguồn: Ok_TracNghiemLop11_Lan1
\begin{ex}%[1D6N1-1]%[ma:712efbff]
	Với $a$ là số thực dương tùy ý, $(\sqrt[3]{a})^{10}: a^2=a^{\tfrac{p}{q}}$ với $p, q  \in \mathbb{Z}$ và $\dfrac{p}{q}$ là phân số tối giản. Giá trị của $p+q$ bằng
	\choice 
		 {$23$}
		{\True  $7$}
		 {$8$}
		{$19$} 
	\loigiai{ 
		 Ta có $(\sqrt[3]{a})^{10}: a^2=a^{\tfrac{4}{3}}$.\\
		Vậy $p+q =  7$.
	} 
\end{ex}

% nguồn: Ok_TracNghiemLop11_Lan1
\begin{ex}%[1D6N1-1]%[ma:fff5900d]
	 Với $a$ là số thực dương tùy ý, tích $a^2\cdot a^{\tfrac{1}{3}}$ bằng
	 \choice
		{\True $a^{\tfrac{7}{3}}$} 
		{$a^{\tfrac{2}{3}}$} 
		{$a^{\tfrac{5}{3}}$} 
		 {$a^{\tfrac{4}{3}}$}
	 \loigiai{
		Ta có $a^2\cdot a^{\tfrac{1}{3}}=  a^{\tfrac{7}{3}}$.
	 }
\end{ex}

% nguồn: Ok_TracNghiemLop11_Lan1
\begin{ex}%[1D6N1-1]%[ma:d65438ac]
	Với $a > 0$, $b > 0$, $\alpha$, $\beta$ là các số thực bất kì, đẳng thức nào sau  đây {\bf{sai}}?
	 \choice
		{$\dfrac{a^{\alpha}}{a^{\beta}}=a^{\alpha-\beta}$} 
		{$a^{\alpha}\cdot  a^{\beta}=a^{\alpha+\beta}$}
		 {\True $\dfrac{a^{\alpha}}{b^{\beta}}=\left(\dfrac{a}{b}\right)^{\alpha-\beta}$}
		{$a^{\alpha}\cdot b^{\alpha}=(ab)^{\alpha}$} 
	\loigiai{ 
		 Đẳng thức sai là $\dfrac{a^{\alpha}}{b^{\beta}}=\left(\dfrac{a}{b}\right)^{\alpha-\beta}$.
	} 
\end{ex}

% nguồn: Ok_TracNghiemLop11_Lan1
\begin{ex}%[1D6N1-1]%[ma:f76be661]
	Cho $a$ là số thực dương khác $1$. Khi đó  $\sqrt[4]{a^{\tfrac{2}{3}}}$ bằng
	\choice 
		{$\sqrt[{3}]{a^2}$} 
		 {$a^{\tfrac{8}{3}}$}
		 {$a^{\tfrac{3}{8}}$}
		{\True $\sqrt[{6}]{a}$} 
	 \loigiai{
		Ta có $\sqrt[4]{a^{\tfrac{2}{3}}}=\left(a^{\tfrac{2}{3}}  \right)^{\tfrac{1}{4}}=a^{\tfrac{2}{3}\cdot\tfrac{1}{4}}=a^{\tfrac{1}{6}}=\sqrt[6]{a}$.
	} 
\end{ex}

% nguồn: Ok_TracNghiemLop11_Lan1
\begin{ex}%[1D6N1-1]%[ma:603e8354]
	Với $a$  là số thực dương tùy ý, $\sqrt[14]{a}$ bằng
	\choice 
		 {\True $a^{\tfrac{1}{14}}$}
		{$a^{\sqrt{14}}$} 
		{$a^{14}$} 
		{$\sqrt{a^{14}}$} 
	 \loigiai{
		Ta có $\sqrt[14]{a} = a^{\tfrac{1}{14}}$. 
	 }
\end{ex}

% nguồn: Ok_TracNghiemLop11_Lan1
\begin{ex}%[1D6N1-1]%[ma:ef6a3c20]
	Số $16$ có bao nhiêu căn bậc $4$? 
	\choice 
		 {\True $2$}
		 {$4$}
		{$1$} 
		{$3$} 
	 \loigiai{
	Vì  $2^4=(-2)^4=16$ nên số $16$ có hai căn bậc $4$.}
\end{ex}

% nguồn: Ok_TracNghiemLop11_Lan1
\begin{ex}%[1D6N1-1]%[ma:c53b2be3]
	Cho các số thực $a$, $b$, $n$, $m$  ($a,b>0$). Khẳng định nào sau đây là đúng?
	\choice 
		{$\left( a^m\right)^n=a^{m+n} $} 
		{$(a+b)^m=a^m+b^m$} 
		 {$\dfrac{a^m}{a^n}=\sqrt{a^m}$}
		{\True  $a^m\cdot a^n=a^{m+n}$}
	\loigiai{ 
		Khẳng định đúng là $a^m\cdot  a^n=a^{m+n}$.
	} 
\end{ex}

% nguồn: Ok_Toan11_TracNghiem_Lan2
\begin{ex}%[1D6H1-1]%[ma:18b59c45]
	Cho biểu thức $T=\dfrac{1}{2^{-x-1}}+3\cdot \sqrt{2}^{2x}-4^{\tfrac{x-1}{2}}$. Khi $2^x=\sqrt{3}$ thì giá trị của biểu  thức $T$ là
	 \choice
		{\True $\dfrac{9\sqrt{3}}{2}$} 
		 {$\dfrac{5\sqrt{3}}{2}$}
		 {$\dfrac{3\sqrt{3}}{2}$}
		{$\dfrac{7\sqrt{3}}{2}$} 
	 \loigiai{
		 Ta có
		 \allowdisplaybreaks
		\begin{eqnarray*}
			T  &= & 2^{x+1}+3\cdot \left(\sqrt{2}^2\right)^x-\left(4^{\tfrac{1}{2}}\right)^{x-1}\\
			 &= & 2\cdot 2^x+3\cdot 2^x-\dfrac{1}{2}\cdot 2^x\\
			&= & \dfrac{9}{2}\cdot  2^x=\dfrac{9}{2}\cdot \sqrt{3} =\dfrac{9\sqrt{3}}{2}.
		\end{eqnarray*}
	 }
\end{ex}

% nguồn: Ok_Toan11_TracNghiem_Lan2
\begin{ex}%[1D6H1-1]%[ma:96038e4c]
	Cho $4^x+4^{-x}=7$. Khi đó biểu thức $P=\dfrac{5-2^x-2^{-x}}{3+2^{x+1}+2^{1-x}}=\dfrac{a}{b}$,  với $\dfrac{a}{b}$ tối giản và $a\in\mathbb{N}$, $b\in\mathbb{N^*}$. Tính tổng $a+b$ có giá trị bằng
	 \choice
		{$8$} 
		 {\True $11$}
		 {$17$}
		{$4$} 
	 \loigiai{
		 Ta có $\left(2^x+2^{-x}\right)^2=4^x+4^{-x}+2\cdot 2^x\cdot 2^{-x}=7+2=9$.\\
		Suy ra  $2^x+2^{-x}=3$.\\
		Do đó $P=\dfrac{5-\left(2^x+2^{-x}\right)}{3+2\cdot \left(2^x+2^{-x}\right)}=\dfrac{5-3}{3+2\cdot  3}=\dfrac{2}{9}$.\\
		Vậy $a=2$ và $b=9$, suy ra $a+b=11$. 
	 }
\end{ex}

% nguồn: Ok_Toan11_TracNghiem_Lan2
\begin{ex}%[1D6H1-1]%[ma:ac5538b5]
	Biết $4^x+4^{-x}=14$, tính giá trị của biểu thức  $P=2^x+2^{-x}$.
	\choice 
		{\True $4$} 
		{$16$} 
		 {$\sqrt{17}$}
		{$\pm  4$}
	 \loigiai{
		 Ta có
		\allowdisplaybreaks 
		\begin{eqnarray*}
			&&4^x+4^{-x}=14  \\
			&\Leftrightarrow& \left(2^x)^2+(2^{-x}\right)^2+2  = 16 \\
			 &\Leftrightarrow& \left(2^x+2^{-x}\right)^2 = 16 \\
			 &\Leftrightarrow&
			\hoac{ 
			 &2^x+2^{-x}=4\\
			 &2^x+2^{-x}=-4.
			}\\ 
			&\Leftrightarrow& 2^x+2^{-x}=4\,  \left(\text{vì } 2^x+2^{-x}>0,\,\forall x\in\mathbb R\right).
		\end{eqnarray*}
		Vậy $P=4$. 
	 }
\end{ex}

% nguồn: Ok_Toan11_TracNghiem_Lan2
\begin{ex}%[1D6H1-1]%[ma:edba3970]
	 Cho biểu thức
	 $T=\dfrac{1}{2^{-x-1}}+3\sqrt{2^{2x}}-4^{\,\tfrac{x-1}{2}}$. Khi $2^x=\sqrt{3}$ thì giá trị của biểu thức $T$ là
	\choice 
		{\True  $\dfrac{9\sqrt{3}}{2}$}
		 {$\dfrac{5\sqrt{3}}{2}$}
		 {$\dfrac{3\sqrt{3}}{2}$}
		 {$\dfrac{7\sqrt{3}}{2}$}
	 \loigiai{
		Ta có 
		\allowdisplaybreaks 
		\begin{align*}
			 T&=\dfrac{1}{2^{-x-1}}+3\sqrt{2^{2x}}-4^{\tfrac{x-1}{2}}\\
			 &=2^{x+1}+3\left(\sqrt{2^2}\right)^x-\left(4^{\tfrac12}\right)^{x-1}\\
			&=2\cdot 2^x+3\cdot  2^x-\dfrac12\cdot 2^x\\
			 &=\dfrac{9}{2}\cdot2^x\\
			&=\dfrac{9\sqrt{3}}{2}. 
		\end{align*}
	 }
\end{ex}

% nguồn: Ok_Toan11_TracNghiem_Lan2
\begin{ex}%[1D6H1-1]%[ma:2b55d9db]
	 Giá trị của biểu thức $P=\dfrac{2^3\cdot 2^{-1}+5^{-3}\cdot 5^4}{10^{-3}:10^{-2}-(0{,}1)^0}$ là
	\choice 
		 {$-9$}
		{\True  $-10$}
		 {$10$}
		 {$9$}
	 \loigiai{
		 Ta có
		 \allowdisplaybreaks
		\begin{eqnarray*}
			P &= &  \dfrac{2^3\cdot 2^{-1}+5^{-3}\cdot 5^4}{10^{-3}:10^{-2}-(0{,}1)^0} =\dfrac{2^{3-1}+5^{-3+4}}{10^{-3+2}-1}\\
			 &= & \dfrac{4+5}{10^{-1}-1}= -10.
		\end{eqnarray*}
	} 
\end{ex}

% nguồn: Ok_Toan11_TracNghiem_Lan2
\begin{ex}%[1D6H1-1]%[ma:3b9a30a2]
	Cho $9^x+9^{-x}=23$. Khi  đó biểu thức
	$A=\dfrac{5+3^x+3^{-x}}{1-3^x-3^{-x}}=\dfrac{a}{b}$ với  $\dfrac{a}{b}$ là phân số tối giản và $a$, $b\in\mathbb Z$. Tích $ab$ bằng
	\choice 
		{\True $-10$} 
		 {$10$}
		 {$-8$}
		{$8$} 
	 \loigiai{
		Ta  có
		\begin{eqnarray*}
			&&9^x+9^{-x}=23 \\ 
			&\Leftrightarrow&  \left(3^x+3^{-x}\right)^2=25\\
			&\Leftrightarrow& 3^x+3^{-x}=5 \left( \text{vì }3^x+3^{-x}>0,\ \forall x\in\mathbb R\right). 
		\end{eqnarray*}
		 Suy ra
		 $A=\dfrac{5+3^x+3^{-x}}{1-3^x-3^{-x}}
		 =\dfrac{5+5}{1-5}
		 =\dfrac{-5}{2}$.\\
		Vậy  $ab=-10$.
	} 
\end{ex}

% nguồn: Ok_Toan11_TracNghiem_Lan2
\begin{ex}%[1D6H1-1]%[ma:fb8c1919]
	Cho $4^x+4^{-x}=7$. Biểu thức $P=\dfrac{5+2^x+2^{-x}}{8-4\cdot 2^x-4\cdot 2^{-x}}$ có giá  trị bằng
	 \choice
		 {$P=\dfrac{3}{2}$}
		{$P=-\dfrac{5}{2}$} 
		{$P=2$} 
		{\True $P=-2$} 
	 \loigiai{
		Ta  có $4^x+4^{-x}=7 \Leftrightarrow \left(2^x+2^{-x}\right)^2=9 \Leftrightarrow 2^x+2^{-x}=3$.\\
		 Suy ra $P=\dfrac{5+2^x+2^{-x}}{8-4\cdot \left(2^x+2^{-x}\right)}=\dfrac{5+3}{8-4\cdot 3}=-2$.
	 }
\end{ex}

% nguồn: Ok_Toan11_TracNghiem_Lan2
\begin{ex}%[1D6H1-1]%[ma:a1b16339]
	Cho $9^x+9^{-x}=23$. Khi đó biểu thức $A=\dfrac{5+3^x+3^{-x}}{1-3^x-3^{-x}}=\dfrac{a}{b}$, với $\dfrac{a}{b}$ là  phân số tối giản và $a$, $b\in\mathbb{Z}$. Tích $a\cdot b$ bằng
	 \choice
		 {\True $-10$}
		 {$10$}
		{$-8$} 
		{$8$} 
	\loigiai{ 
		Ta  có $9^x+9^{-x}=23 \Leftrightarrow \left(3^x+3^{-x}\right)^2=25 \Leftrightarrow 3^x+3^{-x}=5$.\\
		Suy ra  $A=\dfrac{5+3^x+3^{-x}}{1-\left(3^x+3^{-x}\right)}=\dfrac{5+5}{1-5}=\dfrac{-5}{2}$.\\
		Vậy  $a\cdot b=-10$.
	 }
\end{ex}

% nguồn: Ok_TracNghiemLop11_Lan1
\begin{ex}%[1D6H1-1]%[ma:e9b93787]
	 Cho $a$ là số thực dương. Giá trị rút gọn của biểu thức $P=a^{\tfrac{4}{3}} \sqrt{a}$ bằng
	 \choice
		 {\True$a^{\tfrac{11}{6}}$}
		 {$a^{\tfrac{10}{3}}$}
		 {$a^{\tfrac{7}{3}}$}
		{$a^{\tfrac{5}{6}}$} 
	 \loigiai{
	Ta có  $P=a^{\tfrac{4}{3}} \sqrt{a}=a^{\tfrac{4}{3}} \cdot a^{\tfrac{1}{2}}=a^{\tfrac{4}{3}+\tfrac{1}{2}}=a^{\tfrac{11}{6}}$.}
\end{ex}

% nguồn: Ok_TracNghiemLop11_Lan1
\begin{ex}%[1D6H1-1]%[ma:554461eb]
	Biểu diễn biểu thức  $f(x)=\sqrt{x\sqrt[31]{x^{10}}}$ với $x>0$. Khi đó giá trị của $f\left(20^{62}\right)$ bằng
	 \choice
		 {$f\left(20^{62}\right)=1$}
		{$f\left(20^{62}\right)=62^{41}$} 
		{\True $f\left(20^{62}\right)=20^{41}$} 
		{$f\left(20^{62}\right)=20^{62}$} 
	 \loigiai{
		Ta có $f(x)=\sqrt{x\sqrt[31]{x^{10}}}=\sqrt{x\cdot  x^{\tfrac{10}{31}}}=\sqrt{x^{\tfrac{41}{31}}}=x^{\tfrac{41}{62}}$.\\
		 Do đó $f\left(20^{62}\right)=\left(20^{62}\right)^{\tfrac{41}{62}}=20^{41}$.
	} 
\end{ex}

% nguồn: Ok_TracNghiemLop11_Lan1
\begin{ex}%[1D6H1-1]%[ma:f139dcf8]
	Rút gọn biểu thức $A=\dfrac{\sqrt[3]{a^5}\cdot  a^{\tfrac{7}{3}}}{a^4\cdot\sqrt[7]{a^{-2}}}$ với $a > 0$ ta được kết quả $A=a^{\tfrac{m}{n}}$, trong đó $m, n\in\mathbb{N}^*$ và $\dfrac{m}{n}$ là phân số tối giản. Khẳng định nào sau đây đúng?
	\choice 
		 {$3m^2-2n=2$}
		{$m^2+n^2=43$} 
		 {\True $2m^2+n=15$}
		{$m^2+n^2=25$} 
	 \loigiai{
		Ta có $A=\dfrac{\sqrt[3]{a^5}\cdot a^{\tfrac{7}{3}}}{a^4\cdot\sqrt[7]{a^{-2}}}=\dfrac{a^{\tfrac{5}{3}}\cdot a^{\tfrac{7}{3}}}{a^4\cdot  a^{-\tfrac{2}{7}}}=\dfrac{a^{\tfrac{5}{3}+\tfrac{7}{3}}}{a^{4-\tfrac{2}{7}}}=\dfrac{a^4}{a^{\tfrac{26}{7}}}=a^{4-\tfrac{26}{7}}=a^{\tfrac{2}{7}}$.\\
		Suy ra $\left\{\begin{array}{l}
		 m=2\\
		n=7 
	\end{array}\right.\Rightarrow 2m^2+n=2\cdot2^2+7=15$.}
\end{ex}

% nguồn: Ok_TracNghiemLop11_Lan1
\begin{ex}%[1D6H1-1]%[ma:02a88986]
	Rút gọn biểu  thức $P=x^{\frac{1}{6}} \cdot \sqrt[3]{x}$ với $x>0$.
	 \choice
		{$P=x^{\frac{1}{8}}$} 
		 {\True $P=\sqrt{x}$}
		{$P=x^{\frac{2}{9}}$} 
		 {$P=x^2$}
	\loigiai{ 
		 Ta có
		 $$
		P = x^{\frac{1}{6}}  \cdot \sqrt[3]{x}
		= x^{\frac{1}{6}} \cdot  x^{\frac{1}{3}}
		= x^{\frac{1}{2}} 
		=  \sqrt{x}.
		$$ 
	} 
\end{ex}

% nguồn: Ok_TracNghiemLop11_Lan1
\begin{ex}%[1D6H1-1]%[ma:3bc74e3c]
	Cho $a,\,b$ là các số thực  dương và thỏa mãn $a=5^x$, $b=3^x$. Giá trị của biểu thức $P=25^x+15^x+27^x$ bằng
	 \choice
		{$P=a^3+ab+b^2$} 
		 {$P=a^2b^2+ab+b^2$}
		{\True $P=a^2+ab+b^3$} 
		 {$P=a{b^3}+ab+b^2$}
	\loigiai{ 
		Ta  có $P=25^x+15^x+27^x=\left(5^2\right)^x+\left(5.3\right)^x+\left(3^3\right)^x=\left(5^x\right)^2+5^x{3^x}+\left(3^x\right)^3$.\\
		 Với $a=5^x$, $b=3^x$ thì $P=a^2+ab+b^3$.
	} 
\end{ex}

% nguồn: Ok_TracNghiemLop11_Lan1
\begin{ex}%[1D6H1-1]%[ma:1b97592c]
	Biết rằng  $10^\alpha=2$; $10^\beta = 5$. Tính $10^{\alpha+\beta}$+ $10^{\alpha-\beta}$.
	 \choice
		{10} 
		{\True  $\dfrac{52}{5}$}
		 {$\dfrac{2}{5}$}
		{$\dfrac{5}{52}$} 
	 \loigiai{
		 Ta có $10^{\alpha+\beta}=10^\alpha \cdot 10^\beta=2\cdot5=10$. \\
		Mà $10^{\alpha-\beta}=\dfrac{10^\alpha}{  10^\beta}=\dfrac{2}{5}$.\\
		Vậy $10^{\alpha+\beta}$+ $10^{\alpha-\beta}$  = $\dfrac{52}{5}$.
	 }
\end{ex}

% nguồn: Ok_TracNghiemLop11_Lan1
\begin{ex}%[1D6H1-1]%[ma:d115e4d2]
	Cho  hàm số $f(x)=x^{\frac{1}{2}}\cdot\left(x^{\frac{1}{2}}-x^{-\frac{1}{2}}\right)$. Hàm số đã cho có đồ thị không đi qua điểm nào sau đây?
	\choice 
		{\True $(2; 2)$} 
		 {$(1; 0)$}
		 {$(2; 1)$}
		{$(3; 2)$} 
	\loigiai{ 
		 Ta có $f(x)=x^{\frac{1}{2}}\cdot\left(x^{\frac{1}{2}}-x^{-\frac{1}{2}}\right)=x-1$.\\
	Suy ra hàm số  đã cho có đồ thị không đi qua điểm $(2; 2)$.}
\end{ex}

% nguồn: Ok_TracNghiemLop11_Lan1
\begin{ex}%[1D6H1-1]%[ma:e2b6ca79]
	Cho $P=\dfrac{a \sqrt{a} \sqrt[3]{a^2}}{(\sqrt[4]{a})^3}$, ký hiệu $x=\sqrt[12]{a}$. Hãy biễu  diễn $P$ theo $x$.
	 \choice
		{$P=x^{12}$} 
		{$P=x^{10}$} 
		{\True  $P=x^{17}$}
		 {$P=x^{\tfrac{17}{12}}$}
	\loigiai{ 
		Ta có $P=\dfrac{a\sqrt{a} \sqrt[3]{a^2}}{(\sqrt[4]{a})^3}=\dfrac{a^{\tfrac{13}{6}}}{a^{\tfrac{3}{4}}}= a^{\tfrac{17}{12}}$.\\ 
		 Mặt khác $x=\sqrt[12]{a}\Leftrightarrow a = x^{12}$.\\
		Vậy $P  = x^{17}$.
	} 
\end{ex}

% nguồn: Ok_Toan11_TracNghiem_Lan2
\begin{ex}%[1D6C1-1]%[ma:2b968a16]
	Để dự báo dân số của một quốc gia, người ta sử dụng công thức $S=A\mathrm{e}^{nr}$ trong đó $A$ là dân  số của năm lấy làm mốc tính, $S$ là dân số sau $n$ năm, $r$ là tỉ lệ tăng dân số hàng năm. Năm $2017$, dân số Việt nam là $93\,671\,600$ người. Giả sử tỉ lệ tăng dân số hàng năm không đổi là $0{,}81\%$, dự báo dân số Việt nam năm $2035$ là bao nhiêu người?
	\choice 
		{$109\,256\,100$} 
		{\True  $108\,374\,700$}
		{$107\,500\,500$} 
		{$108\,311\,100$} 
	 \loigiai{
		Lấy năm $2017$ làm mốc, ta  có $A=93\,671\,600$; $n=2035-2017=18$.\\
		Suy ra dân số Việt Nam vào  năm $2035$ là
		$$S=93\,671\,600\cdot  \mathrm{e}^{18\cdot 0{,}81\%}\approx 108\,374\,700.$$
	 }
\end{ex}

% nguồn: Ok_Toan11_TracNghiem_Lan2
\begin{ex}%[1D6C1-1]%[ma:4c62bd97]
	 Một người gửi tiết kiệm số tiền $80\,000\,000$ đồng với lãi suất $6{,}9\%$/năm. Biết rằng tiền lãi hàng năm được nhập vào tiền gốc, hỏi sau đúng $5$ năm người đó rút được cả tiền gốc lẫn tiền lãi gần với con số nào sau đây?
	\choice 
		{$105\,370\,000$  đồng}
		{\True $111\,680\,000$ đồng} 
		{$107\,667\,000$  đồng}
		{$116\,570\,000$  đồng}
	 \loigiai{
		Gọi $P_0$ là  số tiền gửi ban đầu, $r$ là lãi suất / năm.\\
		Số tiền gốc và  lãi sau năm thứ nhất là $P_1=P_0+P_0\cdot r=P_0\cdot (1+r)$.\\
		 Số tiền gốc và lãi sau năm thứ hai là $P_2=P_1+P_1\cdot r=P_0\cdot (1+r)^2$.\\
		$\ldots  \ldots \ldots \ldots \ldots$\\
		Số tiền gốc và  lãi người đó rút ra được sau $5$ năm là
		$$P_5=P_0\cdot (1+r)^5=  80\,000\,000\cdot (1+6{,}9\%)^5\approx 111\,680\,000 \text{ (đồng)}.$$
	} 
\end{ex}

% nguồn: Ok_Toan11_TracNghiem_Lan2
\begin{ex}%[1D6C1-1]%[ma:1eb2e866]
	Anh Nam gửi $100$ triệu đồng vào ngân hàng theo thể thức lãi kép, kì hạn là một quý với lãi suất $3\%$ một quý. Sau đúng  $6$ tháng anh Nam gửi thêm $100$ triệu đồng với kì hạn và lãi suất như trước đó. Hỏi sau $1$ năm kể từ lần gửi đầu tiên, số tiền anh Nam nhận được là bao nhiêu?
	\choice 
		 {\True $218{,}64$ triệu đồng}
		{$208{,}25$  triệu đồng}
		{$210{,}45$ triệu  đồng}
		 {$209{,}25$ triệu đồng}
	\loigiai{ 
		Số tiền  anh Nam nhận được sau 6 tháng là
		$$T_1=100\cdot (1+3\%)^2=106{,}09  \text{ (triệu đồng)}.$$
		• Số tiền anh  Nam nhận được sau một năm kể từ lần gửi đầu tiên là
		$$T_2=(106{,}09+100)\cdot (1+3\%)^2\approx 218{,}64 \text{ (triệu  đồng)}.$$
	} 
\end{ex}

% nguồn: Ok_Toan11_TracNghiem_Lan2
\begin{ex}%[1D6C1-1]%[ma:d90815bb]
	Một học sinh $A$ khi $15$ tuổi được hưởng tài sản thừa kế $200\,000\,000$ VNĐ. Số tiền này được bảo quản trong ngân hàng $B$ với kì hạn thanh toán $1$ năm và học sinh $A$ chỉ nhận được số tiền này khi $18$ tuổi. Biết rằng khi $18$ tuổi, số tiền mà học sinh $A$ được nhận sẽ là $231\,525\,000$ VNĐ. Vậy lãi suất kì hạn một năm của ngân hàng $B$  là bao nhiêu?
	 \choice
		 {$8\%/$năm}
		 {$7\%/$năm}
		 {$6\%/$năm}
		{\True $5\%/$năm} 
	\loigiai{ 
		Số  tiền nhận được của gốc và lãi là
		$$200\,000\,000\cdot  (1+r)^3=231\,525\,000 \Leftrightarrow r=5\%/\text{năm}.$$
	 }
\end{ex}

% nguồn: Ok_Toan11_TracNghiem_Lan2
\begin{ex}%[1D6C1-1]%[ma:f36c04c0]
	COVID-$19$ là một loại bệnh viêm đường hô hấp cấp do chủng mới của virus Corona bắt nguồn từ Trung Quốc gây ra với tốc độ truyền bệnh rất nhanh. Giả sử ban đầu có $1$ người bị nhiễm bệnh và cứ sau $1$ ngày sẽ lây sang $4$ người khác. Tất cả những người nhiễm bệnh lại tiếp tục lây sang những  người khác với tốc độ như trên. Hỏi sau $7$ ngày sẽ có tổng cộng bao nhiêu người nhiễm bệnh?
	\choice 
		{$77\,760$ người} 
		{$16\,384$ người} 
		{$62\,500$  người}
		{\True $78\,125$  người}
	 \loigiai{
		 Sau $1$ ngày, tổng số người nhiễm bệnh là $1+4=5$ người.\\
		Sau $2$ ngày, tổng  số người nhiễm bệnh là $(1+4)+(1+4)\cdot 4=(1+4)^2$ người.\\
		Sau $3$ ngày, tổng số người nhiễm bệnh là $(1+4)^2+(1+4)^2\cdot  4=(1+4)^3$ người.\\
		$\ldots \ldots \ldots \ldots  \ldots$\\
		Suy  ra sau $7$ ngày, tổng số người nhiễm bệnh là $(1+4)^7=78\,125$ người.
	} 
\end{ex}

% nguồn: Ok_Toan11_TracNghiem_Lan2
\begin{ex}%[1D6C1-1]%[ma:5e7e858c]
	Ông An gửi $100$ triệu đồng vào một ngân hàng với lãi suất $0{,}8\%$/tháng. Biết rằng nếu không rút tiền ra khỏi ngân hàng thì cứ sau mỗi tháng số tiền lãi sẽ được nhập vào gốc để tính lãi cho tháng tiếp theo và từ tháng thứ hai trở đi, mỗi tháng ông gửi thêm vào tài khoản với số tiền $2$ triệu đồng. Hỏi sau đúng $2$ năm số tiền ông An nhận được cả gốc lẫn lãi là bao nhiêu?  Biết rằng trong suốt thời gian gửi lãi suất không thay đổi và ông An không rút tiền ra.
	 \choice
		 {$169\,871\,000$ đồng}
		{\True $171\,761\,000$ đồng} 
		 {$173\,807\,000$ đồng}
		{$169\,675\,000$  đồng}
	 \loigiai{
		Với $100$ triệu ban  đầu, số tiền cả lãi và gốc thu được sau hai năm là
		$$T_1=100\cdot (1+0{,}8\%)^{24}\cdot 10^6\approx  121\,074\,524 \text{ (đồng)}.$$
		Mỗi tháng tiếp theo gửi $2$ triệu thì tổng số tiền  cả lãi và gốc là
		$$T_2=\dfrac{2}{0{,}008}\cdot\left[(1+0{,}8\%)^{23}-1\right]\cdot(1+0{,}8\%)\cdot10^6=50\,686\,310 \text{  (đồng)}.$$
		Vậy tổng số tiền  là $T=T_1+T_2\approx 171\,761\,000$ đồng.
	 }
\end{ex}

% nguồn: Ok_Toan11_TracNghiem_Lan2
\begin{ex}%[1D6C1-1]%[ma:d529f494]
	Dân số thế giới được ước tính theo công thức $S=A\mathrm{e}^{ni}$, trong đó $A$ là dân số của năm lấy làm mốc, $S$ là dân số sau $n$ năm, $i$ là tỉ lệ tăng dân số hằng năm. Dân số Việt Nam năm $2019$ là $95{,}5$ triệu người, tỉ lệ tăng dân số hằng năm từ $2009$  đến nay là $1{,}14\%$. Hỏi dân số Việt Nam năm $2009$ gần với số nào nhất trong các số sau?
	 \choice
		{$94{,}4$  triệu người}
		 {\True $85{,}2$ triệu người}
		{$86{,}2$  triệu người}
		 {$83{,}9$ triệu người}
	\loigiai{ 
		Ta có  $S=95{,}5$ triệu người, $n=10$ năm, $i=1{,}14\%$.\\
		Suy ra số dân Việt Nam năm $2009$ là  $A=\dfrac{S}{\mathrm{e}^{ni}}=\dfrac{95{,}5}{\mathrm{e}^{10\cdot 1{,}14\%}}\approx 85{,}2$ triệu người.
	 }
\end{ex}

% nguồn: Ok_Toan11_TracNghiem_Lan2
\begin{ex}%[1D6C1-1]%[ma:4fe0b9ac]
	 Ông Chính gửi $200$ triệu đồng vào một ngân hàng với lãi suất $7\%$/năm. Biết rằng nếu không rút tiền ra khỏi ngân hàng thì cứ sau mỗi năm số tiền lãi sẽ được nhập vào gốc để tính lãi cho năm tiếp theo và từ năm thứ $2$ trở đi, mỗi năm ông gửi thêm vào tài khoản với số tiền $20$ triệu đồng. Hỏi sau $18$ năm số tiền ông Chính nhận được cả gốc lẫn lãi là bao nhiêu? Giả định trong suốt thời gian gửi lãi suất không thay đổi và ông Chính không rút tiền ra.
	\choice 
		{$1\,686\,898\,000$ VNĐ} 
		 {$743\,585\,000$ VNĐ}
		{$739\,163\,000$ VNĐ} 
		{\True $1\,335\,967\,000$ VNĐ} 
	\loigiai{ 
		Gọi $a=200$ triệu  đồng; $b=20$ triệu đồng; $\alpha=7\%$.\\
		Số tiền sau $1$ năm là $a\cdot  (1+\alpha)$.\\
		Số tiền sau  $2$ năm là $a\cdot (1+\alpha)^2+b\cdot (1+\alpha)$.\\
		Số tiền sau $3$ năm là $a\cdot  (1+\alpha)^3+b\cdot (1+\alpha)^2+b\cdot (1+\alpha)$.\\
		$\ldots \ldots \ldots \ldots  \ldots$\\
		Vậy số tiền ông Chính nhận sau $18$  năm là
		 \allowdisplaybreaks
		\begin{eqnarray*}
			 & & a\cdot (1+\alpha)^{18}+b\cdot \left[(1+\alpha)^{17}+(1+\alpha)^{17}+(1+\alpha)^{16}+\ldots+(1+\alpha)\right]\\
			&= & a\cdot (1+\alpha)^{18}+b\cdot\left[(1+\alpha)\cdot\dfrac{(1+\alpha)^{17}-1}{\alpha}\right]\\ 
			&\approx & 1\,335\,967\,000 \text{ (VNĐ)}. 
		\end{eqnarray*}
	} 
\end{ex}

% nguồn: Ok_Toan11_TracNghiem_Lan2
\begin{ex}%[1D6C1-1]%[ma:fb0ba37f]
	Vào ngày $15$ hàng tháng ông An đều đến  gửi tiết kiệm tại ngân hàng $SHB$ số tiền $5$ triệu đồng theo hình thức lãi kép với kì hạn một tháng, lãi suất tiết kiệm không đổi trong suốt quá trình gửi là $7{,}2\%$/năm. Hỏi sau đúng $3$ năm kể từ ngày bắt đầu gửi ông An thu được số tiền cả gốc và lãi là bao nhiêu?
	 \choice
		{$195\,251\,000$} 
		{\True $201\,453\,000$} 
		{$195\,252\,000$} 
		{$201\,452\,000$} 
	\loigiai{ 
		Gọi $T_n$ là số tiền  cả gốc lẫn lãi sau $n$ tháng, $A$ là số tiền gốc, $r$ là lãi suất ta có
		\begin{itemize}
			\item Cuối tháng thứ $1$ ông An có số tiền  là $T_1=A(1+r)$.
			 \item Đầu tháng thứ $2$ ông An có số tiền là $T_2=A(1+r)+A$
			\item Cuối tháng thứ $2$ ông An  có số tiền là
			$$T_2=A(1+r)+A+\left[A(1+r)+A\right]r=A(1+r)+A(1+r)^2.$$ 
			$\ldots \ldots \ldots \ldots  \ldots$
			\item Cuối tháng thứ $n$ ông An có số tiền là 
			\allowdisplaybreaks 
			\begin{eqnarray*}
				T_n &= & A(1+r)+A(1+r)^2+\cdots+A(1+r)^n\\ 
				&= &  a\cdot\left[(1+r)+(1+r)^2+\cdots+(1+r)^n\right]\\
				&=  & \dfrac{a(1+r)\left[(1+r)^n-1\right]}{r}.
			\end{eqnarray*}
		\end{itemize}
		Với kì  hạn một tháng, suy ra $3$ năm có $36$ kỳ.\\
		Lãi suất của một năm là $7{,}2\%$, suy ra lãi  suất của $1$ tháng là $\dfrac{7{,}2}{12}\%=0{,}6\%$.\\
		Vậy $T_{36}=\dfrac{5\,000\,000\cdot (1+0{,}6\%)\cdot\left[(1+0{,}6\%)^{36}-1\right]}{0{,}6\%}\approx 201\,453\,000$  (đồng).
	 }
\end{ex}

% nguồn: Ok_Toan11_TracNghiem_Lan2
\begin{ex}%[1D6C1-1]%[ma:274483c6]
	Ngày $01$ tháng $01$ năm $2017$, ông An đem $800$ triệu đồng gửi vào một ngân hàng với lãi suất $0{,}5\%$ một tháng. Từ đó cứ tròn  mỗi tháng, ông đến ngân hàng rút $6$ triệu để chi tiêu cho gia đình. Hỏi đến ngày $01$ tháng $01$ năm $2018$, sau khi rút tiền, số tiền tiết kiệm của ông An còn lại là bao nhiêu? Biết rằng lãi suất trong suốt thời gian ông An gửi không thay đổi.
	 \choice
		 {$800\cdot (1{,}005)^{11}-72$}
		{\True $1\,200-400\cdot  (1{,}005)^{12}$}
		{$800\cdot  (1{,}005)^{12}-72$}
		{$1\,200-400\cdot (1{,}005)^{11}$} 
	\loigiai{ 
		Gửi ngân hàng số tiền là $A$ đồng với lãi suất  $r\%$/tháng và mỗi tháng vào ngày ngân hàng tính lãi, ông An rút ra số tiền là $X$ đồng.\\
		 Sô tiền còn lại sau $n$ tháng đươc tính theo công thức
		 \allowdisplaybreaks
		\begin{eqnarray*}
			 S_n &= & A(1+r)^n-X\cdot\dfrac{(1+r)^n-1}{r}\\
			&= & 800\cdot (1{,}005)^{12}-6\cdot\dfrac{(1{,}005)^{12}-1}{0{,}5\%}\\ 
			&= & 1\,200-400\cdot  (1{,}005)^{12}.
		\end{eqnarray*}
	} 
\end{ex}

% nguồn: Ok_Toan11_TracNghiem_Lan2
\begin{ex}%[1D6C1-1]%[ma:b23e903f]
	Ông Tuấn gửi $100$ triệu vào ngân hàng với hình thức lãi kép, kỳ hạn $1$ năm với lãi suất $8\%$ . Sau $5$ năm ông rút  toàn bộ tiền và dùng một nửa để sửa nhà, số tiền còn lại ông tiếp tục gửi ngân hàng với lãi suất như lần trước. Số tiền lãi ông Tuấn nhận được sau $10$ năm gửi gần nhất với giá trị nào dưới đây?
	\choice 
		{$46{,}933$} 
		 {$34{,}480$}
		{\True $81{,}413$} 
		{$107{,}946$} 
	\loigiai{ 
		Sau $5$ năm đầu, ông Tuấn có số tiền cả gốc và lãi là $T_1=100\cdot  (1+0{,}08)^5$.\\
		Sau khi sửa nhà, số tiền còn lại gửi vào ngân hàng trong $5$ năm thì số tiền cả gốc  và lãi là
		$$T_2=\dfrac{T_1}{2}\cdot (1+0{,}08)^5\approx  107{,}946.$$
		Vậy số tiền lãi trong  $10$ năm là
		$$L=\left(T_1-100\right)+\left(T_2-\dfrac{1}{2}T_1\right)\approx  81{,}413 \text{ (triệu đồng)}.$$
	} 
\end{ex}

% nguồn: Ok_Toan11_TracNghiem_Lan2
\begin{ex}%[1D6C1-1]%[ma:c7de45da]
	Anh An gửi số tiền $58$ triệu đồng vào một ngân hàng theo hình thức lãi kép và ổn định trong $9$ tháng thì lĩnh về được  $61\,758\,000$ đồng. Hỏi lãi suất ngân hàng hàng tháng là bao nhiêu? Biết rằng lãi suất không thay đổi trong thời gian gửi.
	\choice 
		{$0{,}8\%$} 
		{$0{,}6\%$} 
		{\True $0{,}7\%$} 
		 {$0{,}5\%$}
	\loigiai{ 
		Áp dụng công thức $A_n=A_0(1+r)^n$ với $n$ là số kỳ hạn, $A_0$ là số tiền ban  đầu, $A_n$ là số tiền có được sau $n$ kỳ hạn, $r$ là lãi suất.\\
		 Suy ra $A_9=A_0(1+r)^9 \Rightarrow r=\sqrt[9]{\dfrac{A_9}{A_0}}-1=0{,}7\%$.
	 }
\end{ex}

% nguồn: Ok_Toan11_TracNghiem_Lan2
\begin{ex}%[1D6C1-1]%[ma:2548607b]
	Ông An gửi vào ngân hàng $60$ triệu đồng theo hình thức lãi kép. Lãi suất ngân hàng là $8\%$ trên năm. Sau $5$ năm ông  An tiếp tục gửi thêm $60$ triệu đồng nữa. Hỏi sau $10$ năm kể từ lần gửi đầu tiên ông An đến rút toàn bộ tiền gốc và tiền lãi được là bao nhiêu triệu đồng?
	 \choice
		 {$231{,}815$}
		{$197{,}201$} 
		{\True $217{,}695$} 
		{$190{,}271$} 
	\loigiai{ 
		 Số tiền ông An nhận được sau $5$ năm đầu là $A_5=60\cdot (1+8\%)^5$ (triệu đồng).\\
		Số tiền ông An  nhận được sau $10$ năm là
		$$(A_5+60)\cdot  (1+8\%)^5\approx 217{,}695 \text{ (triệu đồng)}.$$
	} 
\end{ex}

% nguồn: Ok_Toan11_TracNghiem_Lan2
\begin{ex}%[1D6C1-1]%[ma:7e23eaa9]
	Ông An gửi tiết kiệm $50$ triệu đồng vào ngân hàng với kỳ hạn $3$ tháng, lãi suất $8{,}4\%$ một năm theo  hình thức lãi kép. Ông gửi được đúng $3$ kỳ hạn thì ngân hàng thay đổi lãi suất, ông gửi tiếp $12$ tháng nữa với kỳ hạn như cũ và lãi suất trong thời gian này là $12\%$ một năm thì ông rút tiền về. Số tiền ông An nhận được cả gốc lẫn lãi là
	\choice 
		{$62\,255\,910$  đồng}
		{\True $59\,895\,767$ đồng} 
		 {$59\,993\,756$ đồng}
		 {$63\,545\,193$ đồng}
	 \loigiai{
		\begin{itemize}
			\item Đợt $I$, ông An gửi số tiền $P_0=50$ triệu đồng, lãi suất $8{,}4\%$ một năm tức là $2{,}1\%$ mỗi  kỳ hạn.\\
			Số  tiền cả gốc và lãi ông thu được sau $3$ kỳ hạn là $P_3=50\,000\,000\cdot (1{,}021)^3$.
			\item  Đợt $II$, do ông không rút ra nên số tiền $P_3$ được xem là số tiền gửi ban đầu của đợt $II$, lãi suất đợt $II$ là $3\%$ mỗi kỳ hạn.\\
			Ông gửi tiếp $12$ tháng (bằng $4$ kỳ hạn) nên số  tiền thu được cuối cùng là
			$$P=P_3\cdot (1{,}03)^4=50\,000\,000\cdot  (1{,}021)^3\cdot (1{,}03)^4 \approx 59\,895\,767 \text{ đồng}.$$
		\end{itemize}
	 }
\end{ex}

% nguồn: Ok_TracNghiemLop11_Lan1
\begin{ex}%[1D6C1-1]%[ma:718c651f]
	Cho $a$, $b$  là các số thực thỏa mãn $3\cdot2^a + 2^b=7\sqrt{2}$ và $5\cdot2^a - 2^b=9\sqrt{2}$. Tính $S=2^{a+b}+2^{a-b}$.
	 \choice
		{$6\sqrt{2}$} 
		 {$2$}
		 {$4$}
		 {\True $6$}
	 \loigiai{
		Ta có  $\left\{
		\begin{aligned}
			3\cdot2^a  + 2^b=7\sqrt{2}\\
			5\cdot2^a  - 2^b=9\sqrt{2}
		\end{aligned}
		\right.$ 
		$\Rightarrow  8\cdot2^a=16\sqrt{2}\Rightarrow 2^a = 2\sqrt{2}$.\\
		Mà $2^b=7\sqrt{2}-3\cdot  2^a\Rightarrow 2^b = 7\sqrt{2}-3\cdot 2\sqrt{2}\Rightarrow 2^b = \sqrt{2}$.\\
		Xét $2^a  \cdot 2^b=2^{a+b}\Rightarrow 2^{a+b}=2\sqrt{2}\cdot \sqrt{2}\Rightarrow 2^{a+b}=4$. \\
		 Mặt khác $\dfrac{2^a}{2^b}=2^{a-b}\Rightarrow 2^{a-b}=\dfrac{2\sqrt{2}}{\sqrt{2}}=2$. \\
		Vậy $S=2^{a+b}+2^{a-b}=4+2=6$. 
		} 
\end{ex}


%%%%%%%%%%%%%%%% ĐÚNG/SAI %%%%%%%%%%%%%%%%

% nguồn: Ok_dungsai_lop11_Lan1
\begin{ex}%[1D6N1-1]%[ma:aa5c9c4f]
 Cho  $a$, $b$ là các số thực dương.
 \choiceTF 
  {$a^m\cdot b^n =  (ab)^{m+n}$}
   {\True $\dfrac{a^m}{a^n}=a^{m-n}$}
   {\True $(a^m)^n=a^{m\cdot n}$}
  {$a^m+a^n=a^{m+n}$} 
 \loigiai{ 
  \begin{itemchoice}
   \itemch Sai.  Vì $a^{m+n}\cdot b^{m+n}=(ab)^{m+n}$.
   \itemch  Đúng.
    \itemch Đúng.
   \itemch Sai. Vì $a^m\cdot a^n=a^{m+n}$. 
  \end{itemchoice}
  }
\end{ex}

% nguồn: Ok_dungsai_lop11_Lan1
\begin{ex}%[1D6N1-1]%[ma:085bebee]
 Cho biểu thức $9^{\tfrac{2}{5}} \cdot  27^{\tfrac{2}{5}}=A$ và $144^{\tfrac{3}{4}}: 9^{\tfrac{3}{4}}=B$, khi đó
  \choiceTF
   {\True $A-B=1$}
  {\True $144^{\tfrac{3}{4}}: 9^{\tfrac{3}{4}}=2^k$ thì $k=3$} 
  {\True $9^{\tfrac{2}{5}} \cdot  27^{\tfrac{2}{5}}=(9 \cdot 27)^{\tfrac{2}{5}}$}
  {$9^{\tfrac{2}{5}} \cdot 27^{\tfrac{2}{5}}=3^k$ thì $k=3$} 
 \loigiai{ 
  \begin{itemchoice}
   \itemch Đúng. Ta có $A-B=9-8=1$. 
   \itemch Đúng.  Ta có $144^{\tfrac{3}{4}}: 9^{\tfrac{3}{4}}=\left(\dfrac{144}{9}\right)^{\tfrac{3}{4}}=16^{\tfrac{3}{4}}=\left(2^4\right)^{\tfrac{3}{4}}=2^3=8=2^3\Rightarrow k=3$.
   \itemch Đúng.  Ta có $9^{\tfrac{2}{5}} \cdot 27^{\tfrac{2}{5}}=(9 \cdot 27)^{\tfrac{2}{5}}$.
   \itemch Sai.  Ta có $9^{\tfrac{2}{5}} \cdot 27^{\tfrac{2}{5}}=(9 \cdot 27)^{\tfrac{2}{5}}=\left(3^2 \cdot 3^3\right)^{\tfrac{2}{5}}=\left(3^5\right)^{\tfrac{2}{5}}=3^2=9=3^2\Rightarrow k =2.$
 \end{itemchoice}}
\end{ex}

% nguồn: Ok_dungsai_lop11_Lan1
\begin{ex}%[1D6H1-1]%[ma:65a841ca]
 Cho biểu thức $9^{\frac{2}{5}}\cdot 27^{\frac{2}{5}}=A$ và $144^{\frac{3}{4}}:9^{\frac{3}{4}}=B$,  khi đó
  \choiceTF
   {\True $9^{\frac{2}{5}}\cdot 27^{\frac{2}{5}}=(9\cdot27)^{\frac{2}{5}}$}
  { $9^{\frac{2}{5}}\cdot 27^{\frac{2}{5}}=3^k$ thì  $k=3$}
   {$144^{\frac{3}{4}}:9^{\frac{3}{4}}=2^k$ thì $k=3$}
  {\True Phép  toán $A-B$ thu được kết quả là một số tự nhiên}
  \loigiai{
  \begin{itemchoice}
   \itemch Đúng. Vì $9^{\frac{2}{5}}\cdot 27^{\frac{2}{5}}=(9\cdot27)^{\frac{2}{5}}$. 
   \itemch Sai.  Vì
    $9^{\frac{2}{5}}\cdot 27^{\frac{2}{5}}=(9\cdot27)^{\frac{2}{5}}=\left(3^2\cdot3^3\right)^{\frac{2}{5}}=(3^5)^{\frac{2}{5}}=3^2=9$.
   \itemch Sai. Vì  $144^{\frac{3}{4}}:9^{\frac{3}{4}}$=$\left(\dfrac{144}{9}\right)^{\frac{3}{4}}=16^{\frac{3}{4}}=(2^4)^\frac{3}{4}=2^3=8$.
   \itemch Đúng. Vì  $A=9^{\frac{2}{5}}\cdot 27^{\frac{2}{5}}=(9\cdot27)^{\frac{2}{5}}=\left(3^2\cdot3^3\right)^{\frac{2}{5}}=(3^5)^{\frac{2}{5}}=3^2=9$.\\
    $B=144^{\frac{3}{4}}:9^{\frac{3}{4}}$=$\left(\dfrac{144}{9}\right)^{\frac{3}{4}}=16^{\frac{3}{4}}=(2^4)^\frac{3}{4}=2^3=8$.\\
    Khi đó $A-B=1$.
  \end{itemchoice}
  }
\end{ex}

% nguồn: Ok_dungsai_lop11_Lan1
\begin{ex}%[1D6H1-1]%[ma:c0a5b832]
 Các câu sau đúng hay  sai?
  \choiceTF
  {\True  $\left[(0{,}2)^{\tfrac{3}{5}}\right]^{-5}=(0{,}2)^{-3}$}
   {\True $\left(5^{\tfrac{2}{3}}\right)^{3}=5^{2}$}
  {$\left(5^{\tfrac{2}{3}}\right)^{-3}+\left[(0{,}2)^{\tfrac{3}{5}}\right]^{-5}=5^m+5^n$ với $m$,  $n$ là các số tự nhiên chẵn}
  {$\left(5^{\tfrac{2}{3}}\right)^{-3}+\left[(0{,}2)^{\tfrac{3}{5}}\right]^{-5}=K$ với $K$ chia  hết cho 4}
  \loigiai{
  \begin{itemchoice}
   \itemch  Đúng. $\left[(0{,}2)^{\tfrac{3}{5}}\right]^{-5}=(0{,}2)^{\tfrac{3}{5}\cdot(-5)}=(0{,}2)^{-3}$
   \itemch Đúng. $\left(5^{\tfrac{2}{3}}\right)^{3}=5^{{\tfrac{2}{3}}\cdot(3)}=5^{2}$. 
    \itemch Sai. $\left(5^{\tfrac{2}{3}}\right)^{3}+\left[(0{,}2)^{\tfrac{3}{5}}\right]^{-5}=5^{2}+(0{,}2)^{-3}=5^{2}+\left(\dfrac{1}{5}\right)^{-3}=5^{2}+5^3\Rightarrow m=2$; $n=3$ là số lẻ.
   \itemch  Sai. $\left(5^{\tfrac{2}{3}}\right)^{3}+\left[(0{,}2)^{\tfrac{3}{5}}\right]^{-5}=5^{2}+(0{,}2)^{-3}=5^{2}+\left(\dfrac{1}{5}\right)^{-3}=5^{2}+5^3=150$.\\
   Mà  $150$ không chia hết cho $4$. Nên $K$ không chia hết cho $4$.
 \end{itemchoice}}
\end{ex}

% nguồn: Ok_dungsai_lop11_Lan1
\begin{ex}%[1D6H1-1]%[ma:2129b0e3]
 Cho $a=\log_2  5$.
  \choiceTF
  {$\log_2 10=\dfrac{1}{a}$+1  }
  {\True  $\log_2 \dfrac{1}{5}=-a$}
  {\True  $4^a=25$}
  {$a^2=5$} 
 \loigiai{ 
  \begin{itemize}
   \item $\log_2 10=\log_2(2\cdot 5)=\log_2 2+\log_2 5=1+a$. 
   \item $\log_2 \dfrac{1}{5}=\log_2  5^{-1}=-\log_2 5=-a$.
   \item $4^{a}=4^{\log_2 5}=2^{2\log_2 5}=2^{\log_2 25}=25$. 
   \item $a^2=\left(\log_2 5\right)^2\ne 25$. 
  \end{itemize}
 } 
\end{ex}

% nguồn: Ok_dungsai_lop11_Lan1
\begin{ex}%[1D6H1-1]%[ma:fafc2a1e]
 Xét đúng sai các  câu sau?
 \choiceTF 
  {\True $81^{-0{,}75}+\left(\dfrac{1}{625}\right)^{-\tfrac{1}{4}}-\left(\dfrac{1}{32}\right)^{\tfrac{3}{5}}=3^m+5-2^n$, với $m+n=0$} 
  {$\left(\dfrac{1}{625}\right)^{-\tfrac{1}{4}}=\left(5^{4}\right)^{-\tfrac{1}{4}}$} 
  {$81^{-0{,}75}+\left(\dfrac{1}{625}\right)^{\tfrac{1}{4}}-\left(\dfrac{1}{32}\right)^{\tfrac{3}{5}}=-\dfrac{a}{b}\left(a, b \in \mathbb{N}^{*}\right)$ và $\dfrac{a}{b}$  là phân số tối giản, khi đó $a-b=52$}
  {\True $81^{-0{,}75}=\left(3^4\right)^{-\tfrac{3}{4}}$} 
 \loigiai{ 
  \begin{itemchoice}
   \itemch Đúng. $81^{-0{,}75}+\left(\dfrac{1}{625}\right)^{-\tfrac{1}{4}}-\left(\dfrac{1}{32}\right)^{\tfrac{3}{5}}=\left(3^4\right)^{-\tfrac{3}{4}}+\left(5^{-4}\right)^{-\tfrac{1}{4}}-\left(2^{-5}\right)^{\tfrac{3}{5}}=3^{-3}+5-2^3\Rightarrow m=-3$; $n=3$.\\ $\Rightarrow m+n=0$. 
   \itemch Sai.  $\left(\dfrac{1}{625}\right)^{-\tfrac{1}{4}}=\left(5^{-4}\right)^{-\tfrac{1}{4}}$.
   \itemch Sai. $81^{-0{,}75}+\left(\dfrac{1}{625}\right)^{-\tfrac{1}{4}}-\left(\dfrac{1}{32}\right)^{\tfrac{3}{5}}=\left(3^4\right)^{-\tfrac{3}{4}}+\left(5^{-4}\right)^{-\tfrac{1}{4}}-\left(2^{-5}\right)^{\tfrac{3}{5}}=3^{-3}+5-2^3=-\dfrac{80}{27}\Rightarrow a =80$; $b=27$.\\ 
   $\Rightarrow a-b=52$. 
   \itemch Đúng. $81^{-0{,}75}=\left(3^4\right)^{-\tfrac{3}{4}}$. 
 \end{itemchoice}}
\end{ex}

% nguồn: Ok_dungsai_lop11_Lan1
\begin{ex}%[1D6V1-1]%[ma:7b86c4d1]
 Tại một xí nghiệp, công thức $P(t)=500\cdot\left(\dfrac{1}{2}\right)^{\tfrac{t}{3}}$ được dùng để tính giá trị  còn lại (tính theo triệu đồng) của một chiếc máy sau thời gian $t$ (tính theo năm) kể từ khi đưa vào sử dụng.
 \choiceTF 
  {\True Giá trị còn lại của máy sau  $3$ năm sử dụng là $250$ triệu đồng}
  {Giá trị còn lại của máy sau  $4$ năm $3$ tháng sử dụng gần bằng $180$ triệu đồng}
  {\True Sau $2$ năm đưa vào sử dưng  thì giá trị của chiếc máy giảm $185$ triệu đồng so với giá trị ban đầu}
   {\True Sau $1$ năm đưa vào sử dụng thì giá trị của chiếc máy giảm $20{,}6\%$ so với giá trị ban đầu của nó}
  \loigiai{
 \begin{itemchoice}
  \itemch \textbf{Đúng}.\\ 
  Giá trị còn lại của máy  sau $3$ năm sử dụng là
   $$
  P(3)=500 \cdot\left(\dfrac{1}{2}\right)^{\tfrac{3}{3}}=250 \text {  (triệu đồng).}
   $$
  \itemch  \textbf{Sai}.\\
  Đổi:  $4$ năm $3$ tháng $=4{,}25$ năm.\\
   Giá trị còn lại của máy sau $4$ năm $3$ tháng sử dụng là
  $$ 
  P(4{,}25)=500\cdot\left(\dfrac{1}{2}\right)^{\tfrac{4{,}25}{3}}  \approx 187 \text { (triệu đồng).}
  $$ 
  \itemch \textbf{Đúng}.\\ 
  Giá trị của  máy lúc ban đầu là $P(0)=500$ (triệu đồng).\\
  Giá trị còn lại của  máy sau $2$ năm đưa vào sử dụng là
   $$
  P(2)=500\cdot\left(\dfrac{1}{2}\right)^{\tfrac{2}{3}} \approx 315 \text { (triệu đồng).  }
  $$ 
  Vậy sau $2$ năm đưa vào sử dụng thì  giá trị của chiếc máy giảm $500-315=185$ triệu đồng so với giá trị ban đầu.
  \itemch \textbf{Đúng}.\\ 
  Giá trị của máy lúc ban  đầu là $P(0)=500$ (triệu đồng).\\
  Giá trị còn  lại của máy sau $1$ năm đưa vào sử dụng là
  $$ 
  P(1)=500 \cdot\left(\dfrac{1}{2}\right)^{\tfrac{1}{3}} \approx 397 \text  { (triệu đồng). }
   $$
  Vậy sau $1$ năm đưa vào sử dụng thì giá trị của chiếc máy giảm $$100 \%-\dfrac{397}{500} \cdot 100 \approx 20{,}6 \%$$ so với giá trị ban đầu  của nó.
 \end{itemchoice}
  }
\end{ex}


%%%%%%%%%%%%%%%% TRẢ LỜI NGẮN %%%%%%%%%%%%%%%%

% nguồn: Ok_TraLoiNgan_Lop11_Lan2
\begin{ex}%[1D6N1-1]%[ma:62393fe9]
	Định luật thứ ba của Kepler về quỹ đạo chuyển động cho biết cách ước tính khoảng thời gian $P$ mà một hành tinh cần để hoàn thành một quỹ đạo quay quanh Mặt Trời. Khoảng thời gian đó được xác định bởi hàm số $P=d^{\tfrac{3}{2}}$ (năm), trong đó $d$ là khoảng cách từ hành tinh đó đến Mặt Trời tính theo đơn vị thiên văn AU. Hỏi Sao Hỏa quay quanh Mặt Trời thì mất bao nhiêu năm Trái Đất?  Biết khoảng cách từ Sao Hỏa đến Mặt Trời là $1{,}52$ AU. (kết quả làm tròn đến hàng đơn vị)
	 \shortans{2}
	\loigiai{ 
		Do Sao Hỏa cách Mặt Trời khoảng cách $d=1{,}52$ (AU) nên khoảng thời gian $P$ mà Sao Hỏa quay quanh Mặt Trời  là
		\begin{align*}
			P=1{,}52^{\tfrac{3}{2}}\approx  2~\text{năm}.
		\end{align*}
	 }
\end{ex}

% nguồn: Ok_TraLoiNgan_Lop11_Lan2
\begin{ex}%[1D6N1-1]%[ma:a2949b5d]
	Năm $2021$, dân số của một quốc gia châu Á là $19$ triệu người. Người ta ước tình rằng dân số của quốc gia này sẽ tăng gấp đôi sau $30$ năm nữa. Khi đó, dân số $A$ của quốc gia đó sau $t$ năm kể từ năm $2021$ được ước tính bằng công thức $A=19\cdot 2^{\tfrac{t}{30}}$. Hỏi với tốc độ tăng dân số như vậy thì sau $20$ năm nữa, dân số  quốc gia này sẽ là bao nhiêu triệu người? (kết quả làm tròn đến hàng đơn vị theo đơn vị tính: triệu người).
	\shortans{30} 
	\loigiai{ 
		Với $t=20$ thì dân số quốc gia này  sau $20$ năm nữa là
		\begin{align*}
			A=19\cdot  2^{\tfrac{20}{30}}\approx 30.
		\end{align*}
	} 
\end{ex}

% nguồn: Ok_TraLoiNgan_Lop11_Lan1
\begin{ex}%[1D6H1-1]%[ma:c6f687ac]
	Biết $2^\alpha+2^{-\alpha}=5$. Giá trị của biểu  thức $P=4^\alpha+4^{-\alpha}$ bằng bao nhiêu?
	\shortans{23} 
	\loigiai{ 
		Ta  có
		\begin{eqnarray*}
			 2^\alpha+2^{-\alpha}=5 &\Leftrightarrow&\left(2^\alpha+2^{-\alpha}\right)^2=25\\
			&\Leftrightarrow& 4^\alpha+4^{-\alpha}+2  \cdot 2^\alpha \cdot 2^{-\alpha}=25 \\
			 &\Leftrightarrow& 4^\alpha+4^{-\alpha}=25-2=23.
		\end{eqnarray*}
	} 
\end{ex}

% nguồn: Ok_TraLoiNgan_Lop11_Lan1
\begin{ex}%[1D6H1-1]%[ma:b708d667]
	Tính  giá trị của biểu thức $T=\dfrac{(2 \sqrt{2}-3)^{2\,024}(2 \sqrt{2}+3)^{2\,025}}{(2 \sqrt{2}-3)^{-1}}$.
	 \shortans{-1}
	 \loigiai{
		Ta có 
		\begin{eqnarray*}
			T&=&\dfrac{(2 \sqrt{2}-3)^{2\,024}(2 \sqrt{2}+3)^{2\,025}}{(2  \sqrt{2}-3)^{-1}}=\dfrac{(2 \sqrt{2}-3)^{2\,024}(2 \sqrt{2}+3)^{2\,024} \cdot(2 \sqrt{2}+3)}{-(2 \sqrt{2}+3)}\\
			&=&-\left[(2  \sqrt{2}-3)(2 \sqrt{2}+3)\right] ^{2\,024}=-(-1)^{2\,024}=-1.
		\end{eqnarray*}
	 }
\end{ex}

% nguồn: Ok_TraLoiNgan_Lop11_Lan1
\begin{ex}%[1D6H1-1]%[ma:4f9ece85]
	Biết $4^x+25^y=10$.  Giá trị của biểu thức $T=\dfrac{x+y}{x y}$ bằng bao nhiêu?
	 \shortans{2}
	 \loigiai{
		Ta có  $4^x=10 \Rightarrow 10^{\tfrac{1}{x}}=4$ và $25^y=10 \Rightarrow 10^{\tfrac{1}{y}}=25$.\\
		Suy  ra
		\begin{eqnarray*}&&10^{\tfrac{1}{x}} \cdot 10^{\tfrac{1}{y}}=4\cdot 25=100\\
			&\Leftrightarrow  &10^{\tfrac{1}{x}+\tfrac{1}{y}}=10^2 \\
			 &\Rightarrow &\dfrac{1}{x}+\dfrac{1}{y}=2\\
			&\Rightarrow &T=\dfrac{x+y}{x y}=2. 
		\end{eqnarray*}
	} 
\end{ex}

% nguồn: Ok_TraLoiNgan_Lop11_Lan1
\begin{ex}%[1D6H1-1]%[ma:d04ab2ad]
	Tính giá trị của  biểu thức $ P=\dfrac{2^3\cdot{2^{-1}}+5^{-3}\cdot{5^4}}{10^{-3}:{10^{-2}}-0{,}1^0}$.
	\shortans{-10} 
	 \loigiai{
	Ta  có $ P=\dfrac{2^3\cdot{2^{-1}}+5^{-3}\cdot{5^4}}{10^{-3}:{10^{-2}}-0,1^0}=\dfrac{2^2+5}{10^{-1}-1}=-10$.}
\end{ex}

% nguồn: Ok_TraLoiNgan_Lop11_Lan1
\begin{ex}%[1D6H1-1]%[ma:5c620f63]
	Cho biểu thức sau $B=\dfrac{a^{\tfrac{4}{3}}\left(a^{-\tfrac{1}{3}}+a^{\tfrac{2}{3}}\right)}{a^{\tfrac{1}{4}}\left(a^{\tfrac{3}{4}}+a^{-\tfrac{1}{4}}\right)}$ với $a>0$. Tính  $B$ khi $a = 100$.
	\shortans{100} 
	 \loigiai{
		 Ta có $B=\dfrac{a^{\tfrac{4}{3}}\left(a^{-\tfrac{1}{3}}+a^{\tfrac{2}{3}}\right)}{a^{\tfrac{1}{4}}\left(a^{\tfrac{3}{4}}+a^{-\tfrac{1}{4}}\right)}=\dfrac{a+a^2}{a+1}=\dfrac{a(a+1)}{a+1}=a$.\\
		Vậy $B  = 100 = 100$.
	} 
\end{ex}

% nguồn: Ok_TraLoiNgan_Lop11_Lan1
\begin{ex}%[1D6H1-1]%[ma:8663329a]
	 Biết $4^x+4^{-x}=23$, tính giá trị biểu thức $ P=2^x+2^{-x}$.
	\shortans{5} 
	\loigiai{ 
		Đặt $ P=2^x+2^{-x}\Rightarrow P  > 0$.\\
		Ta  có
		\begin{eqnarray*}
			 P^2&=&\left(2^x+2^{-x}\right)^2\\
			 &=&4^x+4^{-x}+2\cdot{2^x}\cdot{2^{-x}}\\
			 &=& 23+2=25.
		\end{eqnarray*}
		Do đó $  P=5$.
	} 
\end{ex}

% nguồn: Ok_TraLoiNgan_Lop11_Lan1
\begin{ex}%[1D6H1-1]%[ma:7e9a35d4]
	Cho $4^x+4^{-x}=7$. Tính giá trị của biểu thức  $P=\dfrac{5+2^x+2^{-x}}{8-4\cdot 2^x-4\cdot 2^{-x}}$.
	\shortans{-2} 
	 \loigiai{
		Ta có $4^x+4^{-x}=7 
		 \Leftrightarrow 4^x+2 +4^{-x}=9
		\Leftrightarrow \left( 2^x+2^{-x}\right)^2  =9
		\Rightarrow  2^x+2^{-x}=3$.\\
		 Khi đó $P=\dfrac{5+2^x+2^{-x}}{8-4\cdot 2^x-4\cdot 2^{-x}}
		=\dfrac{5+\left( 2^x+2^{-x}\right) }{8-4\cdot \left( 2^x +2^{-x}\right)  }
		 =\dfrac{5+3}{8-4\cdot 3}=-2$.
	 }
\end{ex}

% nguồn: Ok_TraLoiNgan_Lop11_Lan1
\begin{ex}%[1D6H1-1]%[ma:0f0835b3]
	Biết $9^{\alpha}=\dfrac{1}{2}$. Tính $B=\left(3^{\alpha}+3^{-\alpha}\right)^2-\left(81^{\alpha}-81^{-\alpha}\right)$ (kết quả làm  tròn đến hàng phần trăm).
	 \shortans{8,25}
	\loigiai{ 
		 $$
		\begin{aligned}
			B&=\left(3^\alpha\right)^2+2 \cdot 3^\alpha \cdot  3^{-\alpha}+\left(3^{-\alpha}\right)^2-\left(9^2\right)^\alpha+\left(9^2\right)^{-\alpha} \\
			& =3^{2\alpha}+2 \cdot  3^{\alpha+(-\alpha)}+\left(3^2\right)^{-\alpha}-\left(9^\alpha\right)^2+\left(9^{-\alpha}\right)^2 \\
			 & =9^\alpha+2 \cdot 3^0+9^{-\alpha}-\left(9^\alpha\right)^2+\left(9^\alpha\right)^{-2}=\frac{1}{2}+2+\left(\frac{1}{2}\right)^{-1}-\left(\frac{1}{2}\right)^2+\left(\frac{1}{2}\right)^{-2}\\
			 &=\frac{33}{4}\approx 8{,}25.
		\end{aligned}
		 $$
	 }
\end{ex}

% nguồn: Ok_TraLoiNgan_Lop11_Lan1
\begin{ex}%[1D6H1-1]%[ma:b944735a]
	Biết  rằng $3^x=5$, giá trị của biểu thức $P=81^x+\sqrt[4]{3^x} \cdot \sqrt[4]{27^x}$ bằng bao nhiêu?
	 \shortans{630}
	\loigiai{ 
		Ta có $P=81^x+\sqrt[4]{3^x} \cdot \sqrt[4]{27^x}=\left(3^4\right)^x+\sqrt[4]{3^x  \cdot 27^x}=\left(3^x\right)^4+\sqrt[4]{\left(3^x\right)^4}=5^4+5=630$.
	 }
\end{ex}

% nguồn: Ok_TraLoiNgan_Lop11_Lan1
\begin{ex}%[1D6H1-1]%[ma:e62239d0]
	Cho biểu thức $A=3^{2 x-1}  \cdot\left(\dfrac{1}{3}\right)^{2 x-1}+9^{x+1}$.
	Tính $A$ khi $3^x=2$. 
	\shortans{37} 
	 \loigiai{
		Ta có $A=\left(3 \cdot \dfrac{1}{3}\right)^{2 x-1}+9 \cdot  9^x=1+9\left(3^x\right)^2 = 1+9\left(2\right)^2=37$.
	} 
\end{ex}

% nguồn: Ok_TraLoiNgan_Lop11_Lan2
\begin{ex}%[1D6H1-1]%[ma:f37a8662]
	Biết $\dfrac{\sqrt[3]{2}\cdot \sqrt[4]{2^{-3}}}{\sqrt[6]{2^{-1}}}=2^{\tfrac{m}{n}}$,  với $m\in\mathbb{Z}$, $n\in \mathbb{N}^*$ và nguyên tố cùng nhau. Tính $m\cdot n$.
	\shortans{-4} 
	 \loigiai{
		Ta có 
		\begin{align*}
			\dfrac{\sqrt[3]{2}\cdot \sqrt[4]{2^{-3}}}{\sqrt[6]{2^{-1}}}=\dfrac{2^{\tfrac{1}{3}}\cdot 2^{-\tfrac{3}{4}}}{2^{-\tfrac{1}{6}}}=2^{\tfrac{1}{3}-\tfrac{3}{4}+\tfrac{1}{6}}=2^{-\tfrac{1}{4}} 
		\end{align*}
		Suy ra $m=-1$, $n=4$.\\ 
		 Vậy $m\cdot n=-4$.
	} 
\end{ex}

% nguồn: Ok_TraLoiNgan_Lop11_Lan2
\begin{ex}%[1D6H1-1]%[ma:5d72dcbd]
	Cho $P=\sqrt{2\sqrt{2\sqrt{2\ldots\sqrt{2}}}}$ với  $n$ dấu căn. Tìm $n$, biết $P=2^{1-\tfrac{1}{2^{100}}}$.
	 \shortans{100}
	\loigiai{ 
		Ta có $P=2^{\tfrac{1}{2}}\cdot 2^{\tfrac{1}{2^2}}\cdot  2^{\tfrac{1}{2^3}}\cdots\cdot 2^{\tfrac{1}{2^n}}=2^{\tfrac{1}{2}+\tfrac{1}{2^2}+\tfrac{1}{2^3}+\cdots+\tfrac{1}{2^n}}=2^{\tfrac{1}{2}\cdot \tfrac{\left(\tfrac{1}{2}\right)^n-1}{\tfrac{1}{2}-1}}=2^{1-\tfrac{1}{2^{n}}}$.\\
		Vì $P=2^{1-\tfrac{1}{2^{100}}}$  nên $n=100$.
	} 
\end{ex}

% nguồn: Ok_TraLoiNgan_Lop11_Lan2
\begin{ex}%[1D6H1-1]%[ma:a5792eb9]
	Biết $\sqrt[5]{5\cdot \sqrt[4]{25\cdot \sqrt[3]{5}}} 
	=5^{\tfrac{m}{n}}$ với $m,n\in\mathbb{Z}$ và nguyên tố cùng  nhau. Tính $m+n$.
	 \shortans{79}
	\loigiai{ 
		Ta có 
		\begin{align*}
			\sqrt[5]{5\cdot  \sqrt[4]{25\cdot \sqrt[3]{5}}}
			&=~5^{\tfrac{1}{5}}\cdot  \left(\sqrt[4]{25\cdot 5^{\tfrac{1}{3}}}\right)^{\tfrac{1}{5}}
			=5^{\tfrac{1}{5}}\cdot  \left[\left(5^2\cdot 5^{\tfrac{1}{3}}\right)^{\tfrac{1}{4}}\right]^{\tfrac{1}{5}}\\
			 &=~5^{\tfrac{1}{5}}\cdot \left(5^{\tfrac{7}{12}}\right)^{\tfrac{1}{5}}
			 =5^{\tfrac{1}{5}}\cdot 5^{\tfrac{7}{60}}
			=5^{\tfrac{1}{5}+\tfrac{7}{60}} 
			 =5^{\tfrac{19}{60}}.
		\end{align*}
		Suy  ra $m=19$, $n=60$.\\
		 Vậy $m+n=79$.
	 }
\end{ex}

% nguồn: Ok_TraLoiNgan_Lop11_Lan2
\begin{ex}%[1D6H1-1]%[ma:a521275d]
	Biết $\dfrac{6^{2+\sqrt{5}}\cdot 2^{1-\sqrt{5}}}{3^{3+\sqrt{5}}}=\dfrac{m}{n}$ với $m\in\mathbb{Z}$, $n\in\mathbb{N}^*$ và nguyên  tố cùng nhau. Tính $m+n$.
	\shortans{11} 
	\loigiai{ 
		Ta  có
		\begin{align*}
			\dfrac{6^{2+\sqrt{5}}\cdot 2^{1-\sqrt{5}}}{3^{3+\sqrt{5}}}=\dfrac{2^{2+\sqrt{5}}\cdot 3^{2+\sqrt{5}}\cdot 2^{1-\sqrt{5}}}{3^{3+\sqrt{5}}}=2^{2+\sqrt{5}+1-\sqrt{5}}\cdot 3^{2+\sqrt{5}-3-\sqrt{5}}=2^3\cdot 3^{-1}=\dfrac{8}{3}. 
		\end{align*}
		 Suy ra $m=8$, $n=3$.\\
		 Vậy $m+n=11$.
	 }
\end{ex}

% nguồn: Ok_TraLoiNgan_Lop11_Lan2
\begin{ex}%[1D6H1-1]%[ma:4bebdce1]
	Nếu một khoản tiền gốc $T_0$ được gửi ngân hàng với lãi suất hằng năm $r$, được tính lãi $n$ lần trong một năm, thì tổng số tiền $T_N$ nhận được sau $N$ kì gửi được cho bởi công thức $T_N=T_0\left(1+\dfrac{r}{n}\right)^N$. Hỏi  nếu anh $A$ gửi tiết kiệm số tiền $200$ triệu đồng theo kì hạn $6$ tháng với lãi suất không đổi là $5\%$ /năm thì số tiền thu được của anh $A$ sau $2$ năm khoảng bao nhiêu triệu đồng? (làm tròn đến hàng đơn vị theo đơn vị tính: triệu đồng).
	 \shortans{221}
	\loigiai{ 
		Ta có $T_0=200$, $r=5\%$, $n=2$, $N=4$.\\ 
		Thế  vào công thức ta được số tiền cả vốn lẫn lãi sau $2$ năm là
		\begin{align*}
			T_N=200\cdot \left(1+\dfrac{5\%}{2}\right)^4\approx  221.
		\end{align*}
	} 
\end{ex}

% nguồn: Ok_TraLoiNgan_Lop11_Lan1
\begin{ex}%[1D6V1-1]%[ma:1b52c99d]
	Cho  biểu thức $A=(\sqrt{2}-1)^{2 x}+(3+2 \sqrt{2})^x$. Tính $A$ khi $(\sqrt{2}+1)^x=2$.
	 \shortans{4,25}
	 \loigiai{
		 Ta có $\sqrt{2}-1=\dfrac{(\sqrt{2}-1)(\sqrt{2}+1)}{\sqrt{2}+1}=\dfrac{1}{\sqrt{2}+1}=(\sqrt{2}+1)^{-1}$.\\
		Mà $3+2 \sqrt{2}=(\sqrt{2})^2+2  \sqrt{2}+1=(\sqrt{2}+1)^2$.\\
		 Do đó
		\begin{eqnarray*}
			A&=&\left[(\sqrt{2}+1)^{-1}\right]^{2  x}+\left[(\sqrt{2}+1)^2\right]^x\\
			 &=&(\sqrt{2}+1)^{-2 x}+(\sqrt{2}+1)^{2 x}\\
			 &=&\left[(\sqrt{2}+1)^x\right]^{-2}+\left[(\sqrt{2}+1)^x\right]^2.
		\end{eqnarray*}
		Vậy $A = 4{,}25$. 
	} 
\end{ex}

% nguồn: Ok_TraLoiNgan_Lop11_Lan1
\begin{ex}%[1D6V1-1]%[ma:8c1d18b8]
	Biết $9^\alpha=\dfrac{1}{2}$. Tính $  B=\left(3^\alpha+3^{-\alpha}\right)^2-\left(81^\alpha+81^{-\alpha}\right)$.\\
	\shortans{8,25} 
	 \loigiai{
		Ta  có
		\begin{eqnarray*}
			B&=& \left(3^\alpha\right)^2+2\cdot{3^\alpha}\cdot{3^{-\alpha}}+\left(3^{-\alpha}\right)^2-\left(9^2\right)^\alpha+\left(9^2\right)^{-\alpha}\\ 
			&=& 3^{2\alpha}+2\cdot{3^{\alpha+(-\alpha)}}+\left(3^2\right)^{-\alpha}-\left(9^\alpha\right)^2+\left(9^{-\alpha}\right)^2\\ 
			 &= & 9^\alpha+2\cdot{3^0}+9^{-\alpha}-\left(9^\alpha\right)^2+\left(9^\alpha\right)^{-2}\\
			&= & \dfrac{1}{2}+2+\left(\dfrac{1}{2}\right)^{-1}-\left(\dfrac{1}{2}\right)^2+\left(\dfrac{1}{2}\right)^{-2}\\ 
			 &=&8{,}25.
		\end{eqnarray*}
	} 
\end{ex}

% nguồn: Ok_TraLoiNgan_Lop11_Lan1
\begin{ex}%[1D6V1-1]%[ma:75ac05d1]
	Anh Toàn được tuyển dụng vào một công ty  đầu năm 2013. Công ty trả lương cho anh theo nguyên tắc: Lương khởi điểm anh nhận là $6$ triệu đồng/ tháng và cứ sau $3$ năm công ty lại tăng lương cho anh thêm $25\%$ số lương đang hưởng. Hiện nay (năm 2024) anh đang được hưởng lương là bao nhiêu triệu đồng một tháng? (Kết quả làm tròn đến hàng phần mười)
	 \shortans{11,7}
	 \loigiai{
		Tính từ năm 2013 đến năm  2024, anh Toàn đã được $3$ lần tăng lương.\\
		Lương của anh Toàn sau lần  tăng đầu tiên là $L_1=6\cdot 1{,}25$ triệu.\\
		Lương của anh Toàn sau lần tăng thứ hai  là $L_2=L_1+25\%\cdot L_1=L_1\cdot 1{,}25=6\cdot (1{,}25)^2$ triệu.\\
		Lương của anh Toàn sau lần tăng thứ ba là $L_3=L_2+25\cdot \%L_2=6\cdot (1{,}25)^3\approx  11{,}7 $ triệu.\\
		Vậy  lương của anh Toàn hiện đang hưởng là $11{,}7$ triệu mỗi tháng.
	} 
\end{ex}

% nguồn: Ok_TraLoiNgan_Lop11_Lan1
\begin{ex}%[1D6V1-1]%[ma:14d01f5a]
	Một khu rừng có trữ lượng gỗ là $4\cdot 10^5$\,m$^3$. Biết tốc độ sinh trưởng của các cây lấy gỗ trong khu rừng này là $4 \%$ mỗi năm.  Sau $5$ năm không khai thác, khu rừng sẽ có số mét khối gỗ là $4ab\,c61{,}161$\,m$^3$ với $a$, $b$, $c$ là số tự nhiên. Tính giá trị $a-b+c$.
	\shortans{8} 
	\loigiai{ 
		Nếu trữ lượng gỗ của khu  rừng ban đầu là $A$ thì\\
		Sau năm thứ nhất, lượng gỗ có được là $A+A r=A(1+r), r$ là tốc độ tăng  trưởng/năm.\\
		Sau  năm thứ hai, lượng gỗ có được là $A(1+r)+A(1+r) \cdot r=A(1+r)^2$.\\
		Ta chứng minh được công thức tính lượng gỗ trong khu rừng là $T_n=A(1+r)^n$.\\ 
		Vậy sau  $5$ năm, lượng gỗ trong khu rừng là
		 $$
		T_s=4 \cdot 10^5\left(1+\frac{4}{100}\right)^5=486\,661{,}161\, \mathrm{~m}^3. 
		$$ 
		Do đó, $a=8$; $b=6$; $c=6$. Suy ra $a-b+c=8$. 
	} 
\end{ex}

% nguồn: Ok_TraLoiNgan_Lop11_Lan1
\begin{ex}%[1D6V1-1]%[ma:644ec76a]
	Số lượng vi khuẩn $V$ trong phòng thí nghiệm tính theo công thức $s(t)=s_0\cdot 2^t$ trong đó $s_0$ là số lượng vi khuẩn $V$ lúc đầu, $s(t)$ là số lượng vi khuẩn có trong $t$ phút. Biết sau $3$ phút thì số lượng vi khuẩn $A$ là $625$ nghìn con. Sau $9$ phút thì số lượng vi khuẩn $V$ có dạng $a  \cdot 10^b$ con, với $a$; $b$ là các số tự nhiên, tính giá trị $S=a+b$?
	 \shortans{11}
	\loigiai{ 
		Vì sau $3$ phút thì số lượng vi khuẩn A là $625$ nghìn con nên $625\,000=s_0 \cdot 2^3$.\\ 
		Số lượng vi khuẩn $V$ sau  $9$ phút là
		 $$
		s(t)=\frac{625\,000}{2^3} \cdot 2^9=625\,000  \cdot 2^6=4 \cdot 10^7\text { con}.
		 $$
		Suy ra $a=4$; $b=7$. Từ đó có $S=11$. 
	 }
\end{ex}

% nguồn: Ok_TraLoiNgan_Lop11_Lan2
\begin{ex}%[1D6V1-1]%[ma:7be48599]
	Đồ thị hàm số $y=f(x)$ đối xứng với đồ thị hàm số  $y=\log_a x$ ($0<a\ne 1$) qua điểm $I(2;1)$. Giá trị của biểu thức $f\left(4-a^{-2\,019}\right)$ bằng bao nhiêu?
	\shortans{2021} 
	 \loigiai{
		Xét $y=\log_a x$ có đồ thị là $(C_0)$ và $y=f(x)$ có đồ  thị là $(C)$.\\
		Theo bài ra ta có $(C)$ đối xứng với $(C_0)$ qua  $I(2;1)$.\\
		Giả  sử $M(x_0;y_0)\in (C)$ và $N(x;y)\in (C)$.\\
		Vì  $M$ đối xứng với $N$ qua $I$ nên $I$ là trung điểm của $MN$.\\
		Khi đó $\heva{&x_0+x=4 \\&y_0+y=2}\Leftrightarrow \heva{&x_0=4-x  \\&y_0=2-y.}$\\
		Thay vào phương trình của  $(C_0)$ ta được
		\[2-y=\log_a (4-x)\Leftrightarrow y=2-\log_a (4-x).\] 
		Suy ra  $f\left(4-a^{-2\,019}\right)=2-\log_a \left[4-\left(4-a^{-2\,019}\right)\right]=2-\log_a a^{-2\,019}=2+2\,019=2\,021$.
	 }
\end{ex}


%%%%%%%%%%%%%%%% TỰ LUẬN %%%%%%%%%%%%%%%%

% nguồn: Ok_TuanLuan_Toan11_Lan1
\begin{bt}%[1D6H1-1]%[ma:37482be7]
	Biết  $4^x+4^{-x}=23$, tính giá trị biểu thức $P=2^x+2^{-x}$.
	 \loigiai{
		Đặt $P=2^x+2^{-x} \Rightarrow P >  0$.\\
		Ta có  $P^2=\left(2^x+2^{-x} \right)^2=4^x+4^{-x}+2\cdot 2^x \cdot 2^{-x}=23+2=25$.\\
		Do đó  $P=5$.
	 }
\end{bt}

% nguồn: Ok_TuanLuan_Toan11_Lan1
\begin{bt}%[1D6H1-1]%[ma:e5e079cf]
	Thực hiện các yêu cầu  sau:
	\begin{enumerate}
		\item Cho $4^x+4^{-x}=7$. Tính giá trị biểu thức  $P=\dfrac{5+2^x+2^{-x}}{8-4\cdot2^{-x}-4\cdot2^{-x}}$.
		\item Cho $9^x+9^{-x}=23$. Khi đó biểu thức $A=\dfrac{5+3^x+3^{-x}}{1-3^x-3^{-x}}=\dfrac{a}{b}$ với $\dfrac{a}{b}$ là phân số  tối giản và $a,b\in\mathbb{Z}$. Tính $a\cdot b$.
	\end{enumerate}
	 \loigiai{
		\begin{enumerate}
			\item  Ta có
			$4^x+4^{-x}=7\Leftrightarrow(2^x+2^{-x})^2=9  \Rightarrow 2^x+2^{-x}=3$.\\
			Suy ra $P=\dfrac{5+3}{8-4\cdot2^x-4\cdot2^{-x}}=\dfrac{8}{8-4(2^x+2^{-x})}=\dfrac{8}{8-12}=-2$. 
			\item  Ta có
			$9^x+9^{-x}=23\Leftrightarrow(3^x+3^{-x})^2=25\Rightarrow 3^x+3^{-x}=5$ vì $3^x+3^{-x}>0  ,\forall x\in\mathbb{R}$.\\
			$\Rightarrow A=\dfrac{5+5}{1-5}=\dfrac{10}{-4}=-\dfrac{5}{2}$.\\ 
			Vậy $a\cdot  b=-10$.
		\end{enumerate}
	 }
\end{bt}

% nguồn: Ok_TuanLuan_Toan11_Lan1
\begin{bt}%[1D6H1-1]%[ma:0b4c7b91]
	 Tính giá trị của biểu thức:
	\begin{enumerate}
		 \item $P=\sqrt[5]{-4}\cdot \sqrt[5]{8}$.
		 \item $S=27^{\tfrac{2}{3}}$.
		\item $A=\left(  \dfrac{1}{27}\right) ^{\tfrac{-1}{3}}+\left( \dfrac{1}{16}\right) ^{\tfrac{-3}{4}}$.
	\end{enumerate}
	 \loigiai{
		\begin{enumerate}
			\item Ta có  $P=\sqrt[5]{-4\cdot8}=\sqrt[5]{-32}=-2$.
			\item  Ta có $S=27^{\tfrac{2}{3}}=\sqrt[3]{27^2}=3^2=9$.
			\item $A=\left( \dfrac{1}{27}\right) ^{\tfrac{-1}{3}}+\left( \dfrac{1}{16}\right)  ^{\tfrac{-3}{4}}=27^{\tfrac{1}{3}}+16^{\tfrac{3}{4}}=\sqrt[3]{27}+\sqrt[4]{16^3}=\sqrt[3]{3^3}+\sqrt[4]{2^{12}}=3+2^3=11$.
		\end{enumerate}
	 }
\end{bt}

% nguồn: Ok_TuanLuan_Toan11_Lan1
\begin{bt}%[1D6H1-1]%[ma:322ded3e]
	Số lượng của loại vi khuẩn $A$ trong một phòng thí nghiệm được tính theo công thức $s(t)=s(0)\cdot 2^t$, trong đó $s(0)$ là số lượng vi khuẩn $A$ lúc ban đầu, $s(t)$ là số lượng vi khuẩn $A$ có sau $t$ phút. Biết sau $3$ phút thì số lượng vi khuẩn $A$ là $625$  nghìn con. Hỏi sau $10$ phút thì số lượng vi khuẩn $A$ là bao nhiêu?
	 \loigiai{
		Sau $3$ phút, số lượng vi khuẩn $A$ là $625$  nghìn con nên
		 $$s(3)=s(0)\cdot 2^3\Rightarrow s(0)=\dfrac{s(3)}{2^3}=\dfrac{625000}{8}=78125.$$ Tức là ban đầu có $78125$ con vi khuẩn $A$ trong phòng thí nghiệm.\\
		Sau $10$ phút, số  lượng vi khuẩn là $s(10)=s(0)\cdot 2^{10}=80\cdot 10^6$ (con).
	 }
\end{bt}

% nguồn: Ok_TuanLuan_Toan11_Lan1
\begin{bt}%[1D6H1-1]%[ma:bf497b24]
	Cho $0<a\neq 1$. Tính giá trị của biểu thức  $a^{\log_{\sqrt{a}}4}$.
	 \loigiai{Ta có $a^{\log_{\sqrt{a}}4}=a^{2\log_a4}=a^{\log_a4^2}=4^2=16$.}
\end{bt}

% nguồn: Ok_TuanLuan_Toan11_Lan1
\begin{bt}%[1D6H1-1]%[ma:687d72fb]
	Tính  giá trị của biểu thức:
	\begin{enumerate}
		\item $A=\dfrac{3^4\cdot3^{-2}+2^5\cdot2^{-4}}{2^4\cdot2^3\cdot2^{-2}\cdot3^5\cdot3^{-4}}$. 
		 \item $B=\dfrac{3^5 \cdot\left(\dfrac{1}{3}\right)^3+\left(2^2\right)^{-3}\left(\dfrac{1}{4}\right)^{-4}}{5^{-3} \cdot 25^2+\left(\dfrac{3}{2}\right)^0\left(\dfrac{1}{25}\right)^{-1}}$.
	\end{enumerate}
	\loigiai{ 
		\begin{enumerate}
			 \item $A=\dfrac{3^4\cdot3^{-2}+2^5\cdot2^{-4}}{2^4\cdot2^3\cdot2^{-2}\cdot3^5\cdot3^{-4}}=\dfrac{3^2+2^1}{2^5\cdot3^1}=\dfrac{9+2}{32\cdot3}=\dfrac{11}{96}$.
			\item  $B=\dfrac{3^5 \cdot\left(\dfrac{1}{3}\right)^3+\left(2^2\right)^{-3}\left(\dfrac{1}{4}\right)^{-4}}{5^{-3} \cdot 25^2+\left(\dfrac{3}{2}\right)^0\left(\dfrac{1}{25}\right)^{-1}}=\dfrac{3^2+2^2}{5+5^2}=\dfrac{13}{30}$.
		\end{enumerate}
	 }
\end{bt}

% nguồn: Ok_TuanLuan_Toan11_Lan2
\begin{ex}%[1D6H1-1]%[ma:2cb4b28e]
	Biết $4^x+4^{-x}=23$, tính giá  trị biểu thức $P=2^x+2^{-x}$.
	 \loigiai{
		Đặt $P=2^x+2^{-x} \Rightarrow P > 0$.\\ 
		Ta có $P^2=\left(2^x+2^{-x} \right)^2=4^x+4^{-x}+2\cdot 2^x \cdot  2^{-x}=23+2=25$.\\
		Do  đó $P=5$.
	} 
\end{ex}

% nguồn: Ok_TuanLuan_Toan11_Lan2
\begin{ex}%[1D6H1-1]%[ma:6de0d90b]
	Thực hiện các yêu cầu  sau:
	\begin{enumerate}
		\item Cho $4^x+4^{-x}=7$. Tính giá trị biểu thức $P=\dfrac{5+2^x+2^{-x}}{8-4\cdot2^{-x}-4\cdot2^{-x}}$. 
		\item Cho $9^x+9^{-x}=23$. Khi đó biểu thức  $A=\dfrac{5+3^x+3^{-x}}{1-3^x-3^{-x}}=\dfrac{a}{b}$ với $\dfrac{a}{b}$ là phân số tối giản và $a,b\in\mathbb{Z}$. Tính $a\cdot b$.
	\end{enumerate}
	\loigiai{ 
		\begin{enumerate}
			\item Ta  có
			 $4^x+4^{-x}=7\Leftrightarrow(2^x+2^{-x})^2=9 \Rightarrow 2^x+2^{-x}=3$.\\
			Suy  ra $P=\dfrac{5+3}{8-4\cdot2^x-4\cdot2^{-x}}=\dfrac{8}{8-4(2^x+2^{-x})}=\dfrac{8}{8-12}=-2$.
			\item Ta có 
			$9^x+9^{-x}=23\Leftrightarrow(3^x+3^{-x})^2=25\Rightarrow 3^x+3^{-x}=5$ vì  $3^x+3^{-x}>0 ,\forall x\in\mathbb{R}$.\\
			$\Rightarrow  A=\dfrac{5+5}{1-5}=\dfrac{10}{-4}=-\dfrac{5}{2}$.\\
			Vậy  $a\cdot b=-10$.
		\end{enumerate}
	} 
\end{ex}

% nguồn: Ok_TuanLuan_Toan11_Lan2
\begin{ex}%[1D6H1-1]%[ma:ceec5230]
	Tính giá trị của  biểu thức:
	\begin{enumerate}
		\item $P=\sqrt[5]{-4}\cdot \sqrt[5]{8}$. 
		\item $S=27^{\tfrac{2}{3}}$. 
		\item $A=\left( \dfrac{1}{27}\right)  ^{\tfrac{-1}{3}}+\left( \dfrac{1}{16}\right) ^{\tfrac{-3}{4}}$.
	\end{enumerate}
	 \loigiai{
		\begin{enumerate}
			\item  Ta có $P=\sqrt[5]{-4\cdot8}=\sqrt[5]{-32}=-2$.
			\item Ta có $S=27^{\tfrac{2}{3}}=\sqrt[3]{27^2}=3^2=9$. 
			\item $A=\left( \dfrac{1}{27}\right)  ^{\tfrac{-1}{3}}+\left( \dfrac{1}{16}\right) ^{\tfrac{-3}{4}}=27^{\tfrac{1}{3}}+16^{\tfrac{3}{4}}=\sqrt[3]{27}+\sqrt[4]{16^3}=\sqrt[3]{3^3}+\sqrt[4]{2^{12}}=3+2^3=11$.
		\end{enumerate}
	 }
\end{ex}

% nguồn: Ok_TuanLuan_Toan11_Lan2
\begin{ex}%[1D6H1-1]%[ma:f6d149f9]
	Số lượng của loại vi khuẩn $A$ trong một phòng thí nghiệm được tính theo công thức $s(t)=s(0)\cdot 2^t$, trong đó $s(0)$ là số lượng vi  khuẩn $A$ lúc ban đầu, $s(t)$ là số lượng vi khuẩn $A$ có sau $t$ phút. Biết sau $3$ phút thì số lượng vi khuẩn $A$ là $625$ nghìn con. Hỏi sau $10$ phút thì số lượng vi khuẩn $A$ là bao nhiêu?
	\loigiai{ 
		Sau $3$ phút, số lượng vi khuẩn $A$ là $625$  nghìn con nên
		$$s(3)=s(0)\cdot 2^3\Rightarrow s(0)=\dfrac{s(3)}{2^3}=\dfrac{625000}{8}=78125.$$ Tức là  ban đầu có $78125$ con vi khuẩn $A$ trong phòng thí nghiệm.\\
		Sau $10$ phút, số lượng vi  khuẩn là $s(10)=s(0)\cdot 2^{10}=80\cdot 10^6$ (con).
	} 
\end{ex}

% nguồn: Ok_TuanLuan_Toan11_Lan2
\begin{ex}%[1D6H1-1]%[ma:44195d98]
	Cho  $0<a\neq 1$. Tính giá trị của biểu thức $a^{\log_{\sqrt{a}}4}$.
	\loigiai{Ta  có $a^{\log_{\sqrt{a}}4}=a^{2\log_a4}=a^{\log_a4^2}=4^2=16$.}
\end{ex}

% nguồn: Ok_TuanLuan_Toan11_Lan2
\begin{ex}%[1D6H1-1]%[ma:3fbf6f0d]
	Tính giá trị  của biểu thức:
	\begin{enumerate}
		 \item $A=\dfrac{3^4\cdot3^{-2}+2^5\cdot2^{-4}}{2^4\cdot2^3\cdot2^{-2}\cdot3^5\cdot3^{-4}}$.
		 \item $B=\dfrac{3^5 \cdot\left(\dfrac{1}{3}\right)^3+\left(2^2\right)^{-3}\left(\dfrac{1}{4}\right)^{-4}}{5^{-3} \cdot 25^2+\left(\dfrac{3}{2}\right)^0\left(\dfrac{1}{25}\right)^{-1}}$.
	\end{enumerate}
	\loigiai{ 
		\begin{enumerate}
			\item  $A=\dfrac{3^4\cdot3^{-2}+2^5\cdot2^{-4}}{2^4\cdot2^3\cdot2^{-2}\cdot3^5\cdot3^{-4}}=\dfrac{3^2+2^1}{2^5\cdot3^1}=\dfrac{9+2}{32\cdot3}=\dfrac{11}{96}$.
			\item $B=\dfrac{3^5 \cdot\left(\dfrac{1}{3}\right)^3+\left(2^2\right)^{-3}\left(\dfrac{1}{4}\right)^{-4}}{5^{-3} \cdot  25^2+\left(\dfrac{3}{2}\right)^0\left(\dfrac{1}{25}\right)^{-1}}=\dfrac{3^2+2^2}{5+5^2}=\dfrac{13}{30}$.
		\end{enumerate}
	} 
\end{ex}

% nguồn: Ok_TuanLuan_Toan11_Lan2
\begin{ex}%[1D6H1-1]%[ma:e85ed9f3]
	Giá trị của  $27^{\tfrac{1}{3}}$ bằng
	\loigiai{ 
		 Ta có $27^{\tfrac{1}{3}}=\sqrt[3]{27}=3$.
	 }
\end{ex}

% nguồn: Ok_TuanLuan_Toan11_Lan2
\begin{ex}%[1D6H1-1]%[ma:9f29a2ba]
	Tính giá trị của biểu  thức $P=\sqrt[5]{-4}\cdot\sqrt[5]8$.
	\loigiai{ 
		Ta  có $P=\sqrt[5]{-4}\cdot\sqrt[5]8=\sqrt[5]{\left(-32\right)}=-2$.
	 }
\end{ex}

% nguồn: Ok_TuanLuan_Toan11_Lan2
\begin{ex}%[1D6H1-1]%[ma:9321298c]
	Rút gọn biểu  thức $A=x^{\tfrac{1}{3}}\cdot\sqrt[6]x, x > 0$ ta được
	 \loigiai{
		Ta có  $A=x^{\tfrac{1}{3}}\cdot\sqrt[6]x=x^{\tfrac{1}{2}}=\sqrt{x}$.
	 }
\end{ex}

% nguồn: Ok_TuanLuan_Toan11_Lan2
\begin{ex}%[1D6H1-1]%[ma:3e62a317]
	Cho $a$ là một  số thực dương tùy ý. Viết $a^{\tfrac{2}{3}}\cdot \sqrt{a}$ dưới dạng lũy thừa của $a$ với số mũ hữu tỉ.
	\loigiai{ 
		Ta có  $a^{\tfrac{2}{3}}\cdot\sqrt{a}=a^{\tfrac{2}{3}}\cdot a^{\tfrac{1}{2}}=a^{\tfrac{2}{3}+\tfrac{1}{2}}=a^{\tfrac{7}{6}}$.
	 }
\end{ex}

% nguồn: Ok_TuanLuan_Toan11_Lan2
\begin{ex}%[1D6H1-1]%[ma:3ba8564a]
	Viết biểu thức $P=\sqrt[3]{x\cdot\sqrt[4]x}$ ($x > 0$) dưới dạng luỹ thừa với số  mũ hữu tỷ.
	 \loigiai{
		\[P=\sqrt[3]{x\cdot \sqrt[4]x}=\sqrt[3]{x\cdot x^{\tfrac{1}{4}}}=\left(x^{\tfrac{5}{4}} \right)^{\tfrac{1}{3}}=x^{\tfrac{5}{12}}. 
		\] 
	 }
\end{ex}

% nguồn: Ok_TuanLuan_Toan11_Lan2
\begin{ex}%[1D6H1-1]%[ma:848e839f]
	Rút gọn  biểu thức $P=a\sqrt{a\sqrt[3]a}$, $\left(a > 0\right)$ ta được kết quả là
	 \loigiai{
		Ta có  $P=a\sqrt{a\sqrt[3]a}$ $=a\cdot a^{\tfrac{1}{2}}\cdot a^{\tfrac{1}{6}}=a^{1+\tfrac{1}{2}+\tfrac{1}{6}}=a^{\tfrac{5}{3}}$.
	 }
\end{ex}

% nguồn: Ok_TuanLuan_Toan11_Lan2
\begin{ex}%[1D6H1-1]%[ma:74f32d66]
	Rút gọn biểu thức $Q=a^{\tfrac{7}{3}}:\sqrt[3]a$  với $a > 0$.
	\loigiai{ 
		Ta  có $Q=a^{\tfrac{7}{3}}:\sqrt[3]a=a^{\tfrac{7}{3}}:a^{\tfrac{1}{3}}=a^{\tfrac{7}{3}-\tfrac{1}{3}}=a^2$.
	 }
\end{ex}

% nguồn: Ok_TuanLuan_Toan11_Lan2
\begin{ex}%[1D6H1-1]%[ma:c1ade3bb]
	Rút  gọn biểu thức $P=\dfrac{x^{\tfrac{1}{3}}\cdot \sqrt[6]x}{\sqrt[4]x}$, với $x > 0$.
	 \loigiai{
		Ta  có $P=\dfrac{x^{\tfrac{1}{3}}\cdot \sqrt[6]x}{\sqrt[4]x}=\dfrac{x^{\tfrac{1}{3}}\cdot x^{\tfrac{1}{6}}}{x^{\tfrac{1}{4}}}=x^{\tfrac{1}{3}+\tfrac{1}{6}-\tfrac{1}{4}}=x^{\tfrac{1}{4}}=\sqrt[4]x$.
	 }
\end{ex}

% nguồn: Ok_TuanLuan_Toan11_Lan2
\begin{ex}%[1D6H1-1]%[ma:f6fd807c]
	Rút gọn  biểu thức $P=x^{\tfrac{1}{3}}.\sqrt[6]{\sqrt[3]{x^5}}$ với $x > 0$.
	\loigiai{ 
		 Ta có $P=x^{\tfrac{1}{3}}\cdot x^{\tfrac{5}{18}}=x^{\tfrac{1}{3}+\tfrac{5}{18}}=x^{\tfrac{11}{18}}$.
	 }
\end{ex}

% nguồn: Ok_TuanLuan_Toan11_Lan2
\begin{ex}%[1D6H1-1]%[ma:512fd8e4]
	Cho $a$ là một số thực dương. Viết biểu  thức $P=a^{\tfrac{3}{5}}\cdot\sqrt[3]{a^2}$ dưới dạng lũy thừa với số mũ hữu tỉ.
	 \loigiai{
		Ta có $P=a^{\tfrac{3}{5}}\cdot\sqrt[3]{a^2}=a^{\tfrac{3}{5}}\cdot  a^{\tfrac{2}{3}}=a^{\tfrac{3}{5}+\tfrac{2}{3}}=a^{\tfrac{19}{15}}$.
	 }
\end{ex}

% nguồn: Ok_TuanLuan_Toan11_Lan2
\begin{ex}%[1D6H1-1]%[ma:c0fc55be]
	Cho $a$ là một số thực dương. Giá trị của biểu thức $P=\left(\sqrt{2^a}  \right)^{\tfrac{4}{a}}$ bằng
	 \loigiai{
		Ta có $P=\left(\sqrt{2^a}  \right)^{\tfrac{4}{a}}=\left(2^{\tfrac{a}{2}} \right)^{\tfrac{4}{a}}=2^{^{\tfrac{4}{a}\cdot\tfrac{a}{2}}}=2^2=4$.
	 }
\end{ex}

% nguồn: Ok_TuanLuan_Toan11_Lan1
\begin{bt}%[1D6V1-1]%[ma:6f1062bc]
	Cho số thực $a$ thỏa mãn $0< a\ne 1$. Tính giá trị của biểu thức $T=\log_a \left(\dfrac{a^2 \cdot \sqrt[{3}]{a^2}  \cdot \sqrt[{5}]{a^4}}{\sqrt[{15}]{a^7}} \right)$.
	\loigiai{ 
		Ta có $T=\log_a  \left(\dfrac{a^2 \cdot \sqrt[{3}]{a^2} \cdot \sqrt[{5}]{a^4}}{\sqrt[{15}]{a^7}} \right)=\log_a \frac{a^{2+\frac{2}{3}+\frac{4}{5}}}{a^{\frac{7}{15}}}=\log_a a^{2+\frac{2}{3}+\frac{4}{5}-\frac{7}{15}}=\log_a a^3=3$.
	} 
\end{bt}

% nguồn: Ok_TuanLuan_Toan11_Lan2
\begin{ex}%[1D6V1-1]%[ma:ab6728ab]
	Cho số thực $a$ thỏa mãn $0< a\ne 1$. Tính giá trị của biểu thức $T=\log_a \left(\dfrac{a^2 \cdot \sqrt[{3}]{a^2} \cdot \sqrt[{5}]{a^4}}{\sqrt[{15}]{a^7}} \right)$. 
	\loigiai{ 
		Ta có $T=\log_a \left(\dfrac{a^2 \cdot \sqrt[{3}]{a^2} \cdot \sqrt[{5}]{a^4}}{\sqrt[{15}]{a^7}} \right)=\log_a  \frac{a^{2+\frac{2}{3}+\frac{4}{5}}}{a^{\frac{7}{15}}}=\log_a a^{2+\frac{2}{3}+\frac{4}{5}-\frac{7}{15}}=\log_a a^3=3$.
	 }
\end{ex}

% nguồn: Ok_TuanLuan_Toan11_Lan2
\begin{ex}%[1D6V1-1]%[ma:0afd89df]
	Cho $4^x+4^{-x}=7$. Biểu thức $P=\dfrac{5+2^x+2^{-x}}{8-4\cdot 2^x-4\cdot 2^{-x}}$ có giá trị  bằng
	 \loigiai{
		\[4^x+4^{-x}=7\Leftrightarrow  \left(2^x+2^{-x} \right)^2=9\Leftrightarrow 2^x+2^{-x}=3.
		\] 
		Suy ra $P=\dfrac{5+2^x+2^{-x}}{8-4\cdot  2^x-4\cdot 2^{-x}}=\dfrac{5+3}{8-12}=-2$.
	} 
\end{ex}

% nguồn: Ok_TuanLuan_Toan11_Lan2
\begin{ex}%[1D6V1-1]%[ma:bfcf38f5]
	Cho $a=\dfrac{1}{256}$  và $b=\dfrac{1}{27}$. Tính $A=a^{-\tfrac{3}{4}}+b^{-\tfrac{4}{3}}$.
	 \loigiai{
		Thay $a=\dfrac{1}{256}$, $b=\dfrac{1}{27}$ vào $A=a^{-\tfrac{3}{4}}+b^{-\tfrac{4}{3}}$ ta được 
		\[A=a^{-\tfrac{3}{4}}+b^{-\tfrac{4}{3}}=\left(\dfrac{1}{256} \right)^{-\tfrac{3}{4}}+\left(\dfrac{1}{27} \right)^{-\tfrac{4}{3}}=\left(4^{-4} \right)^{-\tfrac{3}{4}}+\left(3^{-3} \right)^{-\tfrac{4}{3}}=4^3+3^4=145. 
		 \]
	 }
\end{ex}

% nguồn: Ok_TuanLuan_Toan11_Lan2
\begin{ex}%[1D6V1-1]%[ma:787d928d]
	Giá trị của  $A=\left(\dfrac{1}{27} \right)^{-\tfrac{2}{3}}+\left(\dfrac{1}{16} \right)^{-1,25}$ bằng
	\loigiai{ 
		\[A=\left(\dfrac{1}{27} \right)^{-\tfrac{2}{3}}+\left(\dfrac{1}{16}  \right)^{-1,25}=27^{\tfrac{2}{3}}+16^{\tfrac{5}{4}}=\sqrt[3]{27^2}+\sqrt[4]{16^5}=\sqrt[3]{3^6}+\sqrt[4]{2^{20}}=3^2+2^5=41.
		\] 
	} 
\end{ex}

% nguồn: Ok_TuanLuan_Toan11_Lan2
\begin{ex}%[1D6V1-1]%[ma:e0d7759c]
	Cho $4^x+4^{-x}=7$. Khi  đó biểu thức $P=\dfrac{5-2^x-2^{-x}}{3+2^{x+1}+2^{1-x}}=\dfrac{a}{b}$ với $\dfrac{a}{b}$ tối giản và $a\in \mathbb{Z}$, $b\in \mathbb{Z}^{+}$. Tổng $a+b$ có giá trị bằng
	 \loigiai{
		Ta  có $\left(2^x+2^{-x} \right)^2=4^x+4^{-x}+2\cdot 2^x\cdot 2^{-x}$ $=7+2$ $=9$.\\
		Suy ra $2^x+2^{-x}=3$. 
		\[P=\dfrac{5-2^x-2^{-x}}{3+2^{x+1}+2^{1-x}}=\dfrac{5-\left(2^x+2^{-x} \right)}{3+2\left(2^x+2^{-x} \right)}=\dfrac{5-3}{3+2\cdot  3}=\dfrac{2}{9}.
		 \]
		Khi đó $a=2$, $b=9\Rightarrow  a+b=11$.
	 }
\end{ex}

% nguồn: Ok_TuanLuan_Toan11_Lan2
\begin{ex}%[1D6V1-1]%[ma:d19402c5]
	Giá trị của $A=\dfrac{3^5\cdot \left(\dfrac{1}{3} \right)^3+(2^2)^{-3}\cdot \left(\dfrac{1}{4}  \right)^{-4}}{5^{-3}\cdot 25^2+\left(\dfrac{3}{2} \right)^0\cdot\left(\dfrac{1}{25} \right)^{-1}}$ bằng
	 \loigiai{
		\[A=\dfrac{3^5\cdot \left(\dfrac{1}{3}  \right)^3+(2^2)^{-3} \left(\dfrac{1}{4} \right)^{-4}}{5^{-3}\cdot 25^2+\left(\dfrac{3}{2} \right)^0\left(\dfrac{1}{25} \right)^{-1}}=\dfrac{3^2+2^2}{5+5^2}=\dfrac{13}{30}.
		\] 
	} 
\end{ex}

% nguồn: Ok_TuanLuan_Toan11_Lan2
\begin{ex}%[1D6V1-1]%[ma:40f2420f]
	Đơn  giản biểu thức $P=a^{\sqrt{2}}\cdot\left(\dfrac{1}{a} \right)^{\sqrt{2}-1}$ với $a > 0$, được kết quả là
	\loigiai{ 
		 Ta có $P=a^{\sqrt{2}}.\left(\dfrac{1}{a} \right)^{\sqrt{2}-1}=a^{\sqrt{2}}\cdot\left(a^{-1} \right)^{\sqrt{2}-1}=a^{\sqrt{2}}\cdot a^{-\sqrt{2}+1}=a^{\sqrt{2}-\sqrt{2}+1}=a^1=a$.
	} 
\end{ex}

% nguồn: Ok_TuanLuan_Toan11_Lan2
\begin{ex}%[1D6V1-1]%[ma:91bbbbc5]
	Giá trị của $A=\dfrac{3^4\cdot 3^{-2}+2^5\cdot 2^{-4}}{2^4\cdot 2^3-2\cdot 3^5\cdot 3^{-4}}$ bằng 
	 \loigiai{
		 \[A=\dfrac{3^4\cdot 3^{-2}+2^5\cdot 2^{-4}}{2^4\cdot 2^3-2\cdot 3^5\cdot 3^{-4}}=\dfrac{3^2+2}{2^7-2.3}=\dfrac{11}{122}.
		\] 
	 }
\end{ex}

% nguồn: Ok_TuanLuan_Toan11_Lan2
\begin{ex}%[1D6V1-1]%[ma:b35a7a68]
	Cho $9^x+9^{-x}=23$. Khi đó biểu thức $A=\dfrac{5+3^x+3^{-x}}{1-3^x-3^{-x}}=\dfrac{a}{b}$ với  $\dfrac{a}{b}$ là phân số tối giản và $a,b\in \mathbb{Z}$. Tích $a\cdot b$ bằng
	 \loigiai{
		Ta có  $9^x+9^{-x}=23\Leftrightarrow \left(3^x+3^{-x} \right)^2=25$
		$\Leftrightarrow 3^x+3^{-x}=5$ vì $3^x+3^{-x} > 0,\forall x\in  \mathbb{R}$.
		\[\Rightarrow  A=\dfrac{5+3^x+3^{-x}}{1-3^x-3^{-x}}=\dfrac{5+5}{1-5}=\dfrac{-5}{2}.
		\] 
		Suy ra $a=-5$, $b=2$. Vậy $a\cdot b=-10$. 
	 }
\end{ex}

% nguồn: Ok_TuanLuan_Toan11_Lan2
\begin{ex}%[1D6V1-1]%[ma:c4180e03]
	Biết $4^x+4^{-x}=14$, tính giá  trị của biểu thức $P=2^x+2^{-x}$.
	 \loigiai{
		Ta  có
		\allowdisplaybreaks 
		\begin{eqnarray*}
			&&4^x+4^{-x}=14\\ 
			&\Leftrightarrow& \left(2^x \right)^2+\left(2^{-x} \right)^2+2=16\\ &\Leftrightarrow&  \left(2^x+2^{-x} \right)^2=16\\
			&\Leftrightarrow& \hoac{&2^x+2^{-x}=4 \\&  2^x+2^{-x}=-4}\\
			&\Leftrightarrow&  2^x+2^{-x}=4.
		\end{eqnarray*}
		 Vậy $P=4$.
	} 
\end{ex}
